\chapter{个性化排序与响应预估}
\label{chap:rs-ranking}

召回已经把数百万件物品缩小为数百个候选，但登山包、水壶和相机之间仍需作出选择。用户最可能点击的物品，未必带来最高的消费价值；购买概率最高的商品，也可能在没有推荐时就会被购买。评分器必须先回答“分数代表什么”，才能讨论用什么网络计算它。

本章以给定候选后的比较为中心。响应与价值预估描述可能发生的结果，偏好排序学习对象之间的次序，增量排序寻找推荐能改变的结果。确定目标后，字段交互、历史行为和语义内容提供不同的计算依据；这些依据既能由分离模块融合，也能进入联合骨干。日志怎样形成、分数怎样使用与维护，则决定了学到的关系能否支持实际比较。最终展示还涉及对象之间的关系，留到下一章展开。

% Outline: 4.1
\section{预估与排序任务}
\label{sec:rs-ranking-task}

设一次请求的可用信息为$\bm{x}_t$，其中只含时点$t$之前可取得的用户资料、历史$H_{u,t}$及情境；候选集合为$C_t$，候选$i$的属性可以并入评分输入。评分函数$s_\theta(\bm{x}_t,i)$输出一个可比较的数。如果它通过sigmoid形成$p_\theta=\sigma(s_\theta)$，并按点击标签训练，才有了点击概率模型的形式；是否校准还需数据检验。任意单调变换都保持排序次序，却会改变概率解释。

响应必须连同条件和窗口定义。某商品“24小时内购买”的概率，与“曝光后经点击并在7天内购买”的概率不是同一目标。展示位置、页面、价格和库存也会影响响应。排序时尚未确定的位置不能随意填入一个值，再把输出解释成脱离位置的偏好。若任务面向家庭或群组，评分对象还要从个人改为成员集合；第~\ref{sec:rs-agree}、\ref{sec:rs-groupim}~节给出的群组表示可以复用，但成员满意度与群体决策规则仍需分别定义。

预测、排序与决策因此有三种不同的错误。把所有点击概率高估一倍可能不改变次序，却使按期望收益分配资源的决策失准；把前两名交换可能只略微增加平均预测误差，却影响首位展示；准确预测用户喜欢某件物品，也不能证明重复展示十次值得。下文先把这些目标分开，再建立共同评分器。

% Outline: 4.2
\section{预估与排序目标}
\label{sec:rs-ranking-objectives}

选择损失函数之前，应先写出要估计的量及其标签来源。相同的矩阵分解或深层网络，可以拟合显式评分、学习点击排序，也可以在有处理与对照证据时学习推荐增量。架构名称不能替代目标定义。

% Outline: 4.2.1
\subsection{响应与价值预估}
\label{sec:rs-response-value}

令$Y$为某一明确窗口内的二值响应，$V$为响应发生后的价值。把未发生响应时的价值定义为零，则
\begin{equation}
 \mathbb E[YV\mid\bm x,i]
 =\mathbb P(Y=1\mid\bm x,i)\,
   \mathbb E[V\mid Y=1,\bm x,i].
 \label{eq:rs-response-value}
\end{equation}
这个等式来自条件期望，不要求响应与价值独立。教学例中，背包的购买概率为0.04，购买后平均毛利为80元，期望毛利为3.2元；水壶的相应值为0.12和20元，期望毛利为2.4元。购买概率排序选择水壶，期望毛利排序选择背包，二者回答的问题不同。这里的毛利还没有扣除推荐成本，也没有识别增量。

二值响应常用交叉熵$\ell_{\mathrm{BCE}}(y,p)=-y\log p-(1-y)\log(1-p)$；数值价值可用平方损失拟合条件均值，也可按分布和应用选择其他目标。观看时长的零值、截断和长尾应在标签中定义，满意度调查则只覆盖被邀请且实际回答的人，不能把缺失回答填成“不满意”。

% Outline: 4.2.1.1
\subsubsection{多行为联合预估与共享基线}
\label{sec:rs-shared-bottom}

一个视频候选可能同时需要点击、观看和满意度预测。共享底层（shared-bottom）用同一表示网络$f_\theta$处理输入，再由任务头$h_k$输出$\widehat y_k=h_k(f_\theta(\bm x,i))$。它让多个任务共同更新底层参数，而单任务基线为每个目标独立训练一套网络。YouTube多任务排序研究同时使用参与和满意度反馈，恰好说明“容易记录的行为”与“想要了解的体验”并不等价\scite{rs-r16}

设样本$n$在任务$k$的标签可用标记为$m_{n,k}\in\{0,1\}$。一种明确处理缺失的联合目标是
\begin{equation}
 \mathcal L=\sum_{k=1}^{K}\lambda_k
 \frac{\sum_n m_{n,k}\ell_k(y_{n,k},\widehat y_{n,k})}
      {\sum_n m_{n,k}},
 \label{eq:rs-multitask-loss}
\end{equation}
其中只包含当前批次有标签的任务，$\lambda_k$调节训练贡献。按任务分别归一化，可以避免标签多的任务仅凭数量占据全部梯度，但仍需关注损失尺度与采样分布。

例如，一条记录显示点击为1、观看时长为35秒，却没有满意度问卷。若点击预测0.7、时长预测30秒，点击损失为$-\log0.7\approx0.357$；以$\tfrac12(y-\widehat y)^2$为时长损失，则该项为12.5。满意度项被掩码排除。未点击样本的时长是否为零，取决于预测的是“曝光后总观看时长”还是“点击后观看时长”，不能只靠掩码习惯决定。

共享是否有益，应比较相同数据、预算下的各单任务误差。若点击改善而满意度退化，可能存在负迁移，也可能只是标签覆盖或权重变化；应分组检查。式~\eqref{eq:rs-multitask-loss}的权重服务训练，不自动等于最终展示效用的权重。

% Outline: 4.2.1.2
\subsubsection{MMoE：任务门控响应预估}
\label{sec:rs-mmoe}

所有任务固定读取同一个底层表示，可能迫使不同响应共享不合适的信息。\term{多门控混合专家}（multi-gate mixture-of-experts, MMoE）保留一组共享专家，但让每个任务决定如何混合专家输出\scite{rs-r37}对同一候选输入$\bm z$，专家输出为$\bm e_j=f_j(\bm z)$，任务$k$的门控和表示为
\begin{equation}
 \bm g_k=\operatorname{softmax}(\bm W_k\bm z),\qquad
 \bm h_k=\sum_{j=1}^{J}g_{k,j}\bm e_j,\qquad
 \widehat y_k=h_k(\bm h_k).
 \label{eq:rs-mmoe}
\end{equation}
专家由多个任务共享，门控及预测头按任务区分；各响应标签仍使用式~\eqref{eq:rs-multitask-loss}。通用专家网络的容量与路由问题不改变这些标签的含义。

教学例取三个二维专家输出$(1,0),(0,1),(1,1)$。点击门控为$(0.6,0.1,0.3)$，得到$(0.9,0.4)$；满意度门控为$(0.1,0.7,0.2)$，得到$(0.3,0.9)$。若两个教学预测头分别采用$\sigma(h_1-h_2)$与$\sigma(h_2-h_1)$，输出约为0.622和0.646。同样的专家能向两个任务提供不同组合，而任务头仍可进一步变换。

原始MMoE是软混合，通常仍计算全部专家，不能把较小门控权重直接算成省掉的运算。门控相近可以帮助诊断共享情况，却不能证明点击造成满意度或两个行为具有同一因果机制。是否减轻负迁移，仍应逐任务检验并计入专家计算。

% Outline: 4.2.1.3
\subsubsection{PLE：共享与任务专属响应预估}
\label{sec:rs-ple}

MMoE的专家都向各任务开放；若某些信息只对观看有用，另一任务仍可能通过共享专家影响它。\term{渐进分层提取}（progressive layered extraction, PLE）为每个任务保留专属专家，另设共享专家，再逐层交换和提取信息\scite{rs-r38}

在一层中，任务$k$的门控只能混合本任务专家与共享专家，共享门控则可混合各任务及共享专家。得到的任务表示进入下一层的本任务专家，共享表示进入下一层的共享专家。两层以后，任务头读取各自最终表示。较早的共享通路可以传递跨任务知识，较后的任务通路逐渐形成分工；“专属”不意味着这些参数在整个计算图中完全不受其他任务间接影响。

用观看和满意度两个任务作标量教学例。第一层三个专家输出分别为观看2、满意度$-1$、共享1；任务门控在“自身、共享”上分别取$(0.75,0.25)$与$(0.5,0.5)$，产生1.75与0；共享门控取$(0.2,0.3,0.5)$，产生0.6。若第二层的三个专家暂取恒等映射，观看门控$(0.6,0.4)$产生1.29，满意度门控$(0.4,0.6)$产生0.36。再接各自预测头，即可按已有标签联合训练。真实专家当然可以是多层非线性网络。

这个计算展示了门控可见范围，没有证明任务都得到改善。检查PLE时，应同时报告观看和满意度相对单任务、共享底层及MMoE的变化，并固定样本覆盖和预算。多个准确的响应头仍需由第~\ref{chap:rs-constraints}~章的目标与约束决定如何使用。

% Outline: 4.2.1.4
\subsubsection{NMTR：级联行为的联合预估}
\label{sec:rs-nmtr}

浏览、加购与购买还存在业务上规定的行为顺序。\term{神经多任务推荐}（neural multi-task recommendation, NMTR）共享用户、物品嵌入，为各行为保留交互函数，再把前一行为的预测加到后一行为的logit中\scite{rs-r97}以三种行为为例，
\begin{align}
 \widehat y_1&=\sigma(f_1(\bm h_u,\bm v_i)+b_1),\\
 \widehat y_2&=\sigma(f_2(\bm h_u,\bm v_i)+\widehat y_1+b_2),\\
 \widehat y_3&=\sigma(f_3(\bm h_u,\bm v_i)+\widehat y_2+b_3).
 \label{eq:rs-nmtr}
\end{align}
各行为的已观测正例及采样未观测对象形成各自交叉熵，联合更新共享嵌入和行为模块；服务时计算级联输出，以指定目标行为的分数比较候选。

若教学浏览预测为0.6，加购交互logit为$-1$且偏置为零，则加购预测为$\sigma(-0.4)\approx0.401$。若该logit改为1，加购预测会达到$\sigma(1.6)\approx0.832$，高于浏览预测。因此NMTR的级联是预测连接，不是条件概率相乘，也不施加行为概率单调约束。未观测对象仍是带假设的训练负例。

行为顺序应符合当前应用的定义。不能为了凑齐级联而把浏览复制成购买标签；跨设备或其他入口也可能使记录顺序不完整。第~\ref{sec:rs-mbgcn}~节的多行为图利用行为作为输入证据，本节利用多种行为提供监督；广告中的点击后转化概率链路则见第~\ref{chap:rs-advertising}~章。

% Outline: 4.2.2
\subsection{偏好与排序学习}
\label{sec:rs-ranking-learning}

如果目标是把更合适的对象放在前面，不必先得到精确的响应概率。逐对象目标分别拟合每件物品的标签；成对目标直接比较$i$与$j$的分差；列表目标在同一候选列表内分配概率或惩罚次序。例如，以人工相关性$r_i$构造目标分布$q_i$，可最小化$-\sum_{i\in C}q_i\log\operatorname{softmax}(\bm s)_i$。这个归一化只在候选$C$内成立，不等于单件物品被点击的概率。

比较单位改变了训练强调的错误。逐对象损失关心数值偏离，成对损失关心相对关系，列表损失可以强调头部竞争。三者都可以配合不同评分架构，也都要说明候选怎样取得、标签如何观察。列表排序损失仍不直接描述一组物品之间的互补、重复或展示成本。

% Outline: 4.2.2.1
\subsubsection{BPR：隐式反馈成对排序}
\label{sec:rs-bpr}

只有点击、购买等隐式记录时，可以假设用户更偏好已交互物品，而不是采样到的未观测物品。\term{贝叶斯个性化排序}（Bayesian personalized ranking, BPR）把这一假设写成三元组$(u,i,j)$：$i$已交互，$j$未观测\scite{rs-h04}令$\Delta_{u,i,j}=s(u,i)-s(u,j)$，用$\sigma(\Delta_{u,i,j})$建模该偏好关系；结合零均值高斯参数先验，得到负对数后验目标
\begin{equation}
 \mathcal L_{\mathrm{BPR}}
 =-\sum_{(u,i,j)\in D}\log\sigma(\Delta_{u,i,j})
  +\frac{\lambda}{2}\lVert\bm\theta\rVert_2^2.
 \label{eq:rs-bpr}
\end{equation}
对三元组采样作随机梯度更新即可训练。分差越小，正负次序越不确定，$-\log\sigma(\Delta)$越大。

以第~\ref{sec:rs-mf}~节的MF评分为例，$\Delta=\bm h_u^\top(\bm v_i-\bm v_j)$。忽略正则的单步教学更新取$\bm h_u=(0.6,0.2)$、$\bm v_i=(0.4,0.5)$、$\bm v_j=(0.2,0.1)$，得$\Delta=0.2$、损失约0.598。记$g=\sigma(-\Delta)\approx0.4502$，学习率0.1，则同时使用旧向量更新：
\begin{align}
 \bm h'_u&=\bm h_u+0.1g(\bm v_i-\bm v_j)
           \approx(0.6090,0.2180),\\
 \bm v'_i&=\bm v_i+0.1g\bm h_u\approx(0.4270,0.5090),\\
 \bm v'_j&=\bm v_j-0.1g\bm h_u\approx(0.1730,0.0910).
\end{align}
更新后的分差约为0.2458，正物品相对优势增大。加入正则时还应减去对应的参数收缩项。

加权隐式MF拟合二值偏好及置信度，BPR-MF拟合成对关系；相同因子结构不意味着相同目标。BPR也能训练LightGCN等评分器。未观测$j$可能只是没有被展示，故“三元组为真”是学习假设；均匀负采样与热门负采样又赋予不同的比较权重。输出按次序使用，并以相同候选上的排名和截断指标评价，不要求$s$具有概率含义。

% Outline: 4.2.3
\subsection{因果增量与排序学习}
\label{sec:rs-uplift-ranking}

高响应并不等于高增量。设$T\in\{0,1\}$表示是否实施已经定义好的推荐处理，$Y(1),Y(0)$为相应潜在结果。条件平均处理效应
\begin{equation}
 \tau(\bm x,i)=\mathbb E[Y(1)-Y(0)\mid\bm x,i]
 \label{eq:rs-uplift}
\end{equation}
描述同类机会中推荐造成的平均变化。处理可以是“在约定位置展示”，也可以是某个完整推荐方案；改变处理定义就改变了$\tau$。每次机会只能观察一个潜在结果，因此需要对照证据和识别条件，不能把普通响应网络的输出改名为增量。

% Outline: 4.2.3.1
\subsubsection{DLCE：推荐增量的排序学习}
\label{sec:rs-dlce}

教学例中，商品A展示与不展示时的购买率分别为0.9、0.85，商品B为0.6、0.3。按展示后的购买率，A领先；按增量，B的0.3高于A的0.05。\term{DLCE}直接学习这种增量的排序，用推荐处理、购买结果及处理倾向构造因果排名目标\scite{rs-r121}

令$e(\bm x,i)=\mathbb P(T=1\mid\bm x,i)$。在一致性、给定特征后无未观测混杂及$0<e<1$等条件下，观测贡献
\begin{equation}
 \widehat\tau=\frac{TY}{e}-\frac{(1-T)Y}{1-e}
 \label{eq:rs-uplift-ips}
\end{equation}
的条件期望为$\tau$。它可以取较大的正负值，不是单次机会的真实个体效应。因果平均名次把物品名次乘以其效应后平均，期望让正增量对象取得较小名次、负增量对象取得较大名次。将名次展开为候选间的比较，即可形成可训练的排序目标。

DLCE分别保留处理正贡献与对照负贡献，不能直接把带符号的$\widehat\tau$随意乘到普通正权重排序上界：负权重会改变不等式方向。设$\Delta=s(u,i)-s(u,j)$，将对照项的负名次改写为反向比较及与参数无关的常数，得到两种方向的logistic代理：
\begin{equation}
 \ell_{\mathrm{DLCE}}=
 \frac{TY}{\max(e,\chi_1)}\log(1+\exp(-\omega\Delta))
 +\frac{(1-T)Y}{\max(1-e,\chi_0)}
       \log(1+\exp(\omega\Delta)).
 \label{eq:rs-dlce-loss}
\end{equation}
这里$\omega>0$控制平滑程度，$\chi_1,\chi_0$为倾向截断阈值。若需把自然对数形式严格写成0--1比较的上界，可整体除以$\log2$；损失的两种方向不变。处理购买推动$i$上升，对照购买推动$i$下降。原方法从购买记录采样$(u,i)$，再采样比较物品$j$，以MF和正则化随机梯度训练。输出用于增量次序比较，不要求等于校准的$\tau$。

例如处理概率为0.2的一次购买记录，处理项贡献为5；对照概率为0.8的一次购买记录，对照项贡献为$-1.25$。极小概率会放大方差，截断倾向或权重又会引入偏差。若未展示时根本无法追踪购买，就不能把未记录当成$Y(0)=0$。位置纠偏试图恢复响应学习依据，DLCE进一步改变要优化的目标；整体增量实验由第~\ref{chap:rs-evaluation}~章建立。

% Outline: 4.3
\section{候选评分方案}
\label{sec:rs-scoring}

响应、偏好和增量目标都需要一个分数，但输入依据与组织方式可以不同。本节先固定候选，比较字段、历史和语义的处理，再研究它们如何构成完整评分器。观测学习和应用维护是这些方案共同面对的条件，不是另外一种与神经网络并列的模型分类。

% Outline: 4.3.1
\subsection{评分依据与处理方法}
\label{sec:rs-scoring-evidence}

同一背包在工作日通勤与周末登山中可能得到不同响应。字段交互刻画用户、对象及情境的组合；历史评分寻找与候选有关的行为证据；语义评分利用内容及需求的含义。三者可以同时出现在一个模型里。下面比较方法时，默认固定标签窗口、候选范围与请求时可用信息，避免把多读了一类输入误当成计算结构的优势。

% Outline: 4.3.1.1
\subsubsection{基于特征交互的评分}
\label{sec:rs-feature-interaction}

线性模型$s=b+\bm w^\top\bm x$为每个特征赋予固定贡献，却不能单独表达“这类用户在这个页面更喜欢这件物品”。特征交互使一个特征的贡献依赖另一个特征。字段是应用定义的信息类别，如用户ID、物品ID和页面；一个字段可以包含许多类别特征。类别通常查表得到嵌入，连续数值则可以标准化后直接输入，或乘到其嵌入上。

下面的教学输入只激活三个特征，数值均为1，分别属于用户、候选与场景字段。为便于复算，嵌入暂取$\bm e_1=(1,0)$、$\bm e_2=(0,1)$、$\bm e_3=(1,1)$，偏置为$-2$，线性项之和为0.2。真实嵌入由标签学习，并没有把这两个维度预先命名为偏好。比较时需要看清模型交互的是标量、向量、字段对还是分支输出，以及汇总后还保留什么信息。

% Outline: 4.3.1.1.1
\paragraph{FM：二阶特征交互评分}
\label{sec:rs-fm}

为每对稀疏特征各学一个权重，需要大量重复组合。\term{因子分解机}（factorization machine, FM）让特征$a$共享潜在向量$\bm e_a\in\mathbb R^d$，以向量内积产生二阶参数\scite{rs-h05}
\begin{equation}
 s(\bm x)=b+\sum_a w_ax_a+
       \sum_{a<b}\langle\bm e_a,\bm e_b\rangle x_ax_b.
 \label{eq:rs-fm}
\end{equation}
用户特征与不同物品、场景交互时复用同一个$\bm e_a$，使未直接观测的组合也能借助共享参数取得分数。若输入只有用户和物品两个独热字段，交互部分就退化为两侧潜在因子的内积；加入页面、时间等字段后，则还会出现用户—页面、物品—页面等项。

设有$q$个非零特征，直接计算需$O(q^2d)$。展开平方时，每个不同特征对出现两次，因此交互项可以改写为
\begin{equation}
 \frac12\sum_{r=1}^{d}
 \left[\left(\sum_a e_{a,r}x_a\right)^2
       -\sum_a e_{a,r}^2x_a^2\right].
 \label{eq:rs-fm-linear}
\end{equation}
减去自乘项后只剩交叉项，复杂度降为$O(qd)$。在教学输入中，三个内积为0、1、1，故$s=-2+0.2+2=0.2$，点击预测$\sigma(0.2)\approx0.5498$。

CTR训练最小化交叉熵及正则。若该样本点击为1，$\partial\ell/\partial s=p-y\approx-0.4502$，而
$\partial s/\partial\bm e_1=x_1\sum_{a\ne1}x_a\bm e_a=(1,2)$。
忽略正则、学习率取0.1时，$\bm e_1$一次更新为$(1.0450,0.0900)$；其他嵌入、线性项及偏置也由同一误差更新。总类别数为$P$时，嵌入参数约为$Pd$，但从未出现且没有辅助信息的特征仍缺少学习依据，低秩共享不能凭空解决所有冷启动。

% Outline: 4.3.1.1.2
\paragraph{FFM：字段感知交互评分}
\label{sec:rs-ffm}

同一用户与商品交互时需要的信息，可能不同于他与广告位交互时的信息。\term{字段感知因子分解机}（field-aware factorization machine, FFM）让一个特征拥有按对方字段区分的表示\scite{rs-h06}若$f(a)$为特征$a$所属字段，二阶项改为
\begin{equation}
 \sum_{a<b}
 \left\langle\bm e_{a,f(b)},\bm e_{b,f(a)}\right\rangle x_ax_b.
 \label{eq:rs-ffm}
\end{equation}
点击输出和交叉熵可沿用FM，变化在于每一对调用的参数。

例如用户面向物品的向量为$(1,0)$，物品面向用户的向量为$(0.2,0.5)$，贡献为0.2；用户面向页面的向量为$(0,1)$，页面面向用户的向量为$(0.3,0.7)$，贡献为0.7。FM要求用户在这两项中复用同一向量，FFM允许两种关系分开。一次点击误差只更新本样本实际使用的字段对向量及其他评分参数。

若有$F$个字段、$P$个类别和$d$维表示，参数约为$PFd$；一般需要逐对计算，代价约为$O(q^2d)$。它失去了FM式~\eqref{eq:rs-fm-linear}所依赖的单一向量共享。字段划分太细会使参数更稀疏，正则与停止训练时机也更重要。额外自由度是否值得，应在相同输入、标签和资源下检验。

% Outline: 4.3.1.1.3
\paragraph{Wide \& Deep：组合评分}
\label{sec:rs-wide-deep}

已有明确业务组合时，可以保留其直接权重，同时让深层网络学习表示共享。\term{Wide \& Deep}把线性分支与深层分支相加为logit，并联合训练\scite{rs-r31}
\begin{equation}
 p=\sigma\!\left(\bm w_{\mathrm w}^{\top}
       [\bm x;\phi(\bm x)]
       +\bm w_{\mathrm d}^{\top}\bm h_{\mathrm d}+b\right).
 \label{eq:rs-wide-deep}
\end{equation}
$\phi(\bm x)$是显式交叉特征，$\bm h_{\mathrm d}$是类别嵌入及连续输入经过MLP后的表示。显式分支可为“已安装应用A且候选为应用B”设置直接参数，深层分支则在应用、用户与情境表示上共享计算。原论文的应用是Google Play推荐。

假设教学线性分支logit为0.8，深层分支为$-0.3$，偏置为零，则$p=\sigma(0.5)\approx0.6225$；若点击为1，两个分支通过同一$p-y$得到梯度，分别调整交叉权重与嵌入网络。这里是联合训练的同一预测器，不是两个独立概率的平均。

常见的“记忆与泛化”描述对应两种参数使用方式：显式组合可以直接记住已见关系，嵌入共享有机会支持相似但未见的组合。它没有保证新ID已有可靠表示。交叉字典膨胀、字段缺失及两分支时延也需要计入比较。

% Outline: 4.3.1.1.4
\paragraph{DeepFM：低阶与高阶联合评分}
\label{sec:rs-deepfm}

手工构造交叉需要选择哪些组合进入线性分支。\term{DeepFM}改用FM自动形成低阶交互，并让FM与深层分支共享原始特征嵌入\scite{rs-r32}FM分支按式~\eqref{eq:rs-fm}产生分数，深层分支把字段嵌入拼接成向量，输出另一个logit：
\[
 p=\sigma\!\left(s_{\mathrm{FM}}+
       \operatorname{MLP}([\bm e_1x_1;\ldots;\bm e_Fx_F])\right).
\]
同一个点击损失同时更新两条路径及共享嵌入。

在前述三字段例子中，FM分数为0.2。若深层输出0.3，则总概率为$\sigma(0.5)\approx0.6225$。相比Wide \& Deep，低阶交叉不再依赖单独指定每个组合，且两分支共享同一嵌入表。深层分支可以表示复杂非线性关系，但不能仅凭某一隐层就宣布识别了特定的三阶业务作用；其容量、数据需求和时延仍随网络配置改变。

% Outline: 4.3.1.1.5
\paragraph{DCN：有界阶特征交叉评分}
\label{sec:rs-dcn}

直接拼接给MLP时，乘法关系需要网络间接学习。\term{深度交叉网络}（deep \& cross network, DCN）显式让原始输入$\bm z^{(0)}$与当前表示交互\scite{rs-r33}
\begin{equation}
 \bm z^{(\ell+1)}
 =\bm z^{(0)}\bigl(\bm w_\ell^\top\bm z^{(\ell)}\bigr)
   +\bm b_\ell+\bm z^{(\ell)}.
 \label{eq:rs-dcn}
\end{equation}
每层先把当前表示投影成一个标量，再用它缩放原始各维，加上偏置和残差。若只看交叉分支，对原始连续输入的多项式最高阶每层增加一阶，但系数受共享结构约束，并非每个高阶单项式都有独立参数。

教学例取$\bm z^{(0)}=(1,2)$、$\bm w_0=(0.1,0.2)$、偏置为零。第一层标量为0.5，输出$(1.5,3)$；第二层仍取同一教学权重，标量为0.75，输出$(2.25,4.5)$。展开第二层可看到原始输入、二阶项与三阶项共同存在。将交叉表示与并行深层分支拼接，再接sigmoid预测头，即可用点击交叉熵训练。

若输入维度为$D$，原始交叉层的参数和乘加为$O(D)$；它比独立列举所有组合紧凑，也限制了可表示的交互权重。原论文属于ADKDD 2017工作坊。评价时应区分交叉层深度、深层分支容量和新增字段，不把“显式高阶”当作无条件优越性。

% Outline: 4.3.1.1.6
\paragraph{DCN-V2：矩阵交叉与低秩评分}
\label{sec:rs-dcn-v2}

DCN每层用一个标量缩放所有原始维度，交互权重受到限制。\term{DCN-V2}把向量参数推广为矩阵，使不同输出维度读取不同的输入组合\scite{rs-r34}
\begin{equation}
 \bm z^{(\ell+1)}
 =\bm z^{(0)}\odot(\bm W_\ell\bm z^{(\ell)}+\bm b_\ell)
  +\bm z^{(\ell)}.
 \label{eq:rs-dcn-v2}
\end{equation}
注意偏置位于逐元素乘法内部，与式~\eqref{eq:rs-dcn}的原始写法不同。取$\bm z^{(0)}=(1,2)$及
$\bm W_0=\left(\begin{smallmatrix}0.1&0.2\\0.3&0.4\end{smallmatrix}\right)$，
偏置为零，则矩阵乘积为$(0.5,1.1)$，一层输出为$(1.5,4.2)$，第二维不再被同一个0.5缩放。

完整矩阵需要$O(D^2)$参数与计算；近似$\bm W=\bm U\bm V^\top$、秩为$r$时可降为$O(Dr)$。论文还研究低秩专家混合，在低维空间加入非线性和门控。这里的专家用于交叉变换，与MMoE按响应任务路由专家的职责不同。

交叉网络可与MLP并行后拼接，也可先交叉再堆叠MLP，最终用相同响应损失学习。比较完整矩阵、低秩及混合方案，要共同记录误差、参数和候选批次时延。该工作正式发表于WWW 2021，2020年作者稿的日期不应替代会议年份。

% Outline: 4.3.1.1.7
\paragraph{xDeepFM：向量级显式交互评分}
\label{sec:rs-xdeepfm}

交叉还可以把一个字段的整条嵌入作为单位。\term{xDeepFM}中的压缩交互网络（compressed interaction network, CIN）让上一层向量与原始字段向量逐元素相乘，再在线性组合中压缩成新的向量\scite{rs-r35}令$\bm X^{(0)}\in\mathbb R^{F\times d}$为字段矩阵，上一层有$H_{\ell-1}$行，则第$h$个输出向量为
\begin{equation}
 \bm X^{(\ell)}_{h,:}
 =\sum_{a=1}^{H_{\ell-1}}\sum_{b=1}^{F}
   \bm W^{(\ell,h)}_{a,b}
   \bigl(\bm X^{(\ell-1)}_{a,:}\odot\bm X^{(0)}_{b,:}\bigr).
 \label{eq:rs-cin}
\end{equation}
同一字段对的权重跨嵌入维度共享，区别于对每个标量单独设交叉系数。各层输出沿嵌入维求和，汇总后与线性分支、深层分支共同产生点击logit。

在三字段例子中，若第一层仅给$(1,3)$和$(2,3)$两对权重1，其余为0，则输出为$(1,0)+(0,1)=(1,1)$，维度汇总为2。下一层再与原始字段相乘，可以进一步形成向量组合。该例只展示一个压缩通道，真实模型可学习多个通道并接CTR交叉熵。计算随$H_\ell H_{\ell-1}Fd$增加，压缩宽度是容量与代价的选择；它不是简单把FM叠得更深，也不消除其他分支的数据需求。

% Outline: 4.3.1.1.8
\paragraph{AutoInt：注意力特征交互评分}
\label{sec:rs-autoint}

固定交叉权重不随当前样本调整读取关系。\term{AutoInt}把各类别、数值字段变成同维向量，让多头自注意力在字段之间交换信息，再由预测头输出点击概率\scite{rs-r36}一个头计算
\[
 \alpha_{a,b}=
 \frac{\exp(\langle\bm W_Q\bm e_a,\bm W_K\bm e_b\rangle)}
      {\sum_c\exp(\langle\bm W_Q\bm e_a,\bm W_K\bm e_c\rangle)},
 \qquad
 \widetilde{\bm e}_a=\sum_b\alpha_{a,b}\bm W_V\bm e_b.
\]
多头结果拼接，加上投影残差并经过激活，得到下一层字段表示；堆叠以后再接响应头并联合学习。这里写的是原模型的内积注意形式，不必强行添加其他Transformer变体的缩放。

教学例暂令投影为恒等，用户字段$(1,0)$对三个字段的logit为$(1,0,1)$，权重约为$(0.422,0.155,0.422)$，聚合为$(0.845,0.578)$。残差再加回用户字段后，原信息与其他字段信息共同进入后层。字段两两注意通常带来$O(F^2d)$量级的交互计算，投影另需预算。

AutoInt的注意对象是同一候选样本内的字段；DIN的注意对象是与候选匹配的历史行为。两者可以结合。注意权重有助于检查模型读取了什么，但不等于字段的因果作用，更不能据此推断用户真正作选择的心理过程。

% Outline: 4.3.1.1.9
\paragraph{RankMixer：异质特征混合评分}
\label{sec:rs-rankmixer}

扩大交互网络时，既要让字段交换信息，也要保留不同字段的独立变换能力。\term{RankMixer}将用户、候选及已经加工好的历史表示分组为等维单元，通过多头重排混合交换子向量，再由每个单元独立的前馈网络变换\scite{rs-r39}混合本身主要改变数据排列，增加容量的参数主要位于各单元的前馈计算及输入、输出投影。

一个简化重排例子有两个四维单元$(a_1,a_2,a_3,a_4)$与$(b_1,b_2,b_3,b_4)$。按两头切分并交换后，可以形成$(a_1,a_2,b_1,b_2)$及$(a_3,a_4,b_3,b_4)$，各自进入独立MLP。于是一个单元的后续输出同时依赖两个原字段。实际模型按设定的头数和形状执行重排，层间还保留残差及归一化；汇总后由任务头拟合点击等响应。

稀疏专家变体还需区分总参数与每个样本激活的参数。应在固定历史编码器、字段和标签下比较其与DCN-V2、AutoInt的交互代价。RankMixer仍可依赖独立序列模块先加工历史，并未单凭字段混合统一全部行为建模。文献按CIKM 2025会议记录使用，作者稿模板中的占位年份与DOI不构成发表证据。

表~\ref{tab:rs-feature-models}把前述计算放在共同坐标下。表中均可接点击概率头及交叉熵，区别在于交互和共享；这一共同输出不意味着原始输入、容量或训练预算天然相同。
\begin{table*}[htbp]
 \centering\small
 \begin{tabularx}{\linewidth}{@{}lXXXX@{}}
 \toprule
 方法 & 字段与交互 & 参数共享 & 汇总及输出 & 代价与适用条件\\
 \midrule
 FM & 稀疏特征的两两内积 & 每特征一个向量 & 二阶和加线性项 & 稀疏线性计算；低秩假设\\
 FFM & 对方字段参与选向量 & 按特征和对方字段 & 二阶和加线性项 & 更多表项；需字段对数据\\
 Wide \& Deep & 显式交叉及嵌入MLP & 组合权重与表示分工 & 两分支logit相加 & 交叉字典与深层计算\\
 DeepFM & FM及嵌入MLP & 两分支共用嵌入 & 两分支logit相加 & 深层增量成本；联合梯度\\
 DCN & 原输入乘当前投影 & 向量参数约束系数 & 交叉与深层融合 & 交叉层线性代价\\
 DCN-V2 & 矩阵或低秩交叉 & 矩阵、低秩或专家 & 并行或堆叠响应头 & 秩、专家与延迟取舍\\
 xDeepFM & CIN向量级乘积 & 跨维共享组合权重 & 层级汇总加其他分支 & 字段数和压缩宽度\\
 AutoInt & 字段间多头注意 & 共享投影和残差 & 保留字段后预测 & 字段两两读取\\
 RankMixer & 分组单元重排混合 & 单元前馈分别参数化 & 汇总后任务预测 & 激活计算与重排开销\\
 \bottomrule
 \end{tabularx}
 \caption{特征交互评分的共同接口。比较性能时还须固定输入信息、标签窗口及计算预算；不同共享限制没有脱离数据条件的统一优劣次序。}
 \label{tab:rs-feature-models}
\end{table*}

% Outline: 4.3.1.1.10
\paragraph{PNN：乘积层的特征交互评分}
\label{sec:rs-pnn}

FM在内积之后直接求和，深层网络则可以继续利用乘积结果。\term{乘积神经网络}（product-based neural network, PNN）在字段嵌入与MLP之间加入乘积层，组合线性信号和乘积信号后预测响应\scite{rs-r46}内积版本把字段对压成标量，外积版本保留不同潜在维度之间的组合；原论文用分解及汇总结构降低直接构造全部乘积的成本。

对教学用户$\bm e_1=(1,0)$与对象$\bm e_2=(0,1)$，内积为0，而外积为
$\left(\begin{smallmatrix}0&1\\0&0\end{smallmatrix}\right)$。
内积丢弃了“第一维与第二维共同出现”的信息，外积仍保留它。若将三个字段的两两内积$(0,1,1)$送入教学隐单元，权重取$(1,0.5,-0.5)$且偏置为0.2，则ReLU输出0.2；以它直接作为logit得到概率0.5498，点击正例损失约0.598。实际PNN还并入线性信号并使用多层网络。

逐对内积约需$O(F^2d)$，逐对完整外积约需$O(F^2d^2)$；论文降本实现改变了参数化与汇总方式，比较时须指明版本。后续非线性响应头使它一般不能直接转成普通单向量点积索引召回，服务于给定候选评分时应按候选数计算预算。

% Outline: 4.3.1.1.11
\paragraph{NFM：交互池化的非线性评分}
\label{sec:rs-nfm}

FM把潜在维度也一并相加，后续网络便看不到不同维度的组合模式。\term{神经因子分解机}（neural factorization machine, NFM）把二阶交互先汇总成向量，再进行非线性变换\scite{rs-r47}双交互池化为
\begin{equation}
 \bm b(\bm x)=\sum_{a<b}x_a\bm e_a\odot x_b\bm e_b
 =\frac12\left[
 \left(\sum_a x_a\bm e_a\right)^{\odot2}
 -\sum_a(x_a\bm e_a)^{\odot2}\right].
 \label{eq:rs-nfm}
\end{equation}
右侧同样只需$O(qd)$。预测由偏置、线性项及$\bm b$的网络输出组成；无隐层且输出权重全为1时，恢复FM的二阶和。

三字段例子的交互向量为$(1,1)$。若教学输出权重为$(1,-1)$且没有隐层，则交互输出为0，总预测为$-1.8$；若此时回归标签为1，以半平方误差计损失为3.92。该人为权重用来说明两个维度可以承担不同作用，不代表论文训练结果。原论文以稀疏回归形式评估上下文预测与标签推荐；迁移CTR时，须再定义sigmoid头、点击标签及交叉熵。

NFM保留了潜在维度，却在池化时丢失了具体交互对的身份。若需要判断“哪一对字段更应参与”，就需要在汇总之前作选择。

% Outline: 4.3.1.1.12
\Needspace{90mm}
\paragraph{AFM：注意权重的二阶交互评分}
\label{sec:rs-afm}

\term{注意因子分解机}（attentional factorization machine, AFM）为每个二阶交互向量$\bm b_{a,b}=x_ax_b(\bm e_a\odot\bm e_b)$计算权重，再加权汇总\scite{rs-r48}注意网络可写成
$a'_{a,b}=\bm h^\top\operatorname{ReLU}(\bm W\bm b_{a,b}+\bm c)$，
在本样本所有有效特征对上softmax，得$a_{a,b}$。输出投影$\bm p^\top\sum_{a<b}a_{a,b}\bm b_{a,b}$与偏置、线性项合成预测。

教学三对向量为$(0,0),(1,0),(0,1)$。若注意权重为$(0.1,0.7,0.2)$，汇总为$(0.7,0.2)$；均匀权重则为$(1/3,1/3)$。以$\bm p=(1,-1)$投影，二者分别贡献0.5与0，说明样本相关权重能改变潜在维度的合成。注意网络、嵌入及预测头由任务损失共同学习。

原论文采用稀疏回归，CTR迁移仍需明确更换输出及损失。与NFM相比，逐对注意增加了$O(q^2)$个交互对象的处理，不能直接沿用其全部线性化优势。高权重表示当前模型重视这一交互，不能证明该字段对造成了点击。

% Outline: 4.3.1.1.13
\paragraph{FiBiNET：字段选择与双线性交互评分}
\label{sec:rs-fibinet}

AFM在交互之后赋权，\term{FiBiNET}还在交互之前调整字段本身。它用压缩激励网络从当前字段嵌入提取摘要，生成字段权重$a_f$，得到重加权嵌入$\widetilde{\bm e}_f=a_f\bm e_f$；原始和重加权两组嵌入分别计算双线性交互，再由浅层或深层预测头输出CTR\scite{rs-r49}

采用列向量记法，一对交互可写成$\bm b_{a,b}=(\bm W\bm e_a)\odot\bm e_b$。矩阵$\bm W$可以全局共享、按一个字段共享，或为每个字段对单独设置。它在逐元素相乘前混合潜在维度，区别于FM直接内积。假设用户权重为0.5、场景权重为1.5，$\bm e_1=(1,0)$、$\bm e_3=(1,1)$，且$\bm W=\left(\begin{smallmatrix}1&0\\1&1\end{smallmatrix}\right)$，原始交互为$(1,1)$，重加权交互为$(0.75,0.75)$。两者拼接后仍可由预测头分别利用。

字段门控、双线性矩阵与点击头通过交叉熵联合学习。矩阵参数从$O(d^2)$增加到按字段的$O(Fd^2)$或按字段对的$O(F^2d^2)$，对应更细的关系和更高的数据需求。门控诊断、交互诊断及实际效果识别属于不同证据，不能互相替代。

% Outline: 4.3.1.1.14
\paragraph{Fi-GNN：字段图的交互评分}
\label{sec:rs-fi-gnn}

若希望每个字段经过多轮交换后再参与预测，可以把同一候选样本的字段构成图。\term{Fi-GNN}以字段为节点，以字段间的学习边权传播消息，再用GRU和初始状态残差更新节点\scite{rs-r50}这张图只有当前样本的用户、物品、页面等节点，并不是目录中的用户—物品交互图。

节点初始化来自字段表示，边权由两端初始表示计算。每个节点具有输入和输出变换，边变换由发送端与接收端的矩阵组合，从而不必为每条边独立存一整套矩阵。以列向量约定写一轮消息：
\[
 \bm m_a=\sum_{b\ne a}\alpha_{b,a}
       \bm W^{\mathrm{in}}_a\bm W^{\mathrm{out}}_b\bm h_b,\qquad
 \bm h'_a=\operatorname{GRU}(\bm h_a,\bm m_a)+\bm h_a^{(0)}.
\]
若教学变换为恒等，物品及场景传给用户的权重为0.6、0.4，则消息为$0.6(0,1)+0.4(1,1)=(0.4,1)$。GRU决定保留多少旧状态及新消息，初始残差帮助保留字段身份。

若干轮后，节点预测与节点权重聚合成整体响应，用点击交叉熵训练。字段数为$F$时，全连接图仍有$O(F^2)$条边，每轮消息和状态更新均有成本；矩阵共享减少的是参数存储，并不使所有字段对消失。本节采用2019会议版机制，后续作者稿的修改应与它区分。节点、边的可视化说明模型的计算依赖，不构成因果解释。

% Outline: 4.3.1.1.15
\paragraph{DLRM：离散与连续特征的联合评分}
\label{sec:rs-dlrm}

价格、时段统计等连续输入不宜总被当成类别。\term{深度学习推荐模型}（deep learning recommendation model, DLRM）让离散字段查嵌入表，连续输入经过底部MLP映射为同维稠密向量，再计算两两点积，并将点积及稠密向量交给顶部MLP\scite{rs-r51}

例如用户、物品两条嵌入为$(1,0),(0,1)$，连续特征经底部MLP得到$(1,1)$。三个向量形成的非重复两两内积为$(0,1,1)$；顶部输入拼接为$(1,1,0,1,1)$，维度为$2+\binom32=5$。顶部MLP最后输出一个logit，sigmoid后以点击交叉熵训练，梯度同时到达稀疏表、底部及顶部网络。若稠密向量和$F$个稀疏向量都为$d$维，顶部输入维度为$d+\binom{F+1}{2}$。

FM直接累加各对交互，DLRM保留每对内积供顶部网络分别变换。显式操作虽然是二阶，经过非线性顶部后，整体函数不受二阶多项式限制。参数常集中于类别表，稠密计算集中于MLP与交互；两者并行实现的细节属于系统卷。改变采样或输入处理时仍需重新检查概率校准。

% Outline: 4.3.1.1.16
\paragraph{MaskNet：实例门控的特征评分}
\label{sec:rs-masknet}

同一用户换了候选，某些表示维度的重要程度也会变化。\term{MaskNet}从当前原始嵌入生成逐维掩码，对嵌入或隐层作乘法缩放，并与前馈网络、层归一化组成MaskBlock\scite{rs-r52}掩码由压缩再扩展的网络产生，通常不是0或1的硬选择；每个块可从同一原始输入取得自己的掩码。

假设某层归一化表示为$(1,-1,0.5,-0.5)$。候选A对应的教学掩码为$(1,0.2,0.5,1)$，乘积为$(1,-0.2,0.25,-0.5)$；候选B掩码为$(0.2,1,1,0.5)$，乘积为$(0.2,-1,0.5,-0.25)$。即使用户相同，候选参与生成掩码后也会改变送入前馈层的表示。接预测头后即可按点击标签计算交叉熵并反向更新。

串联MaskNet将一个块的输出交给下一个块，并联MaskNet让多个块处理共同输入，再拼接送入网络。FiBiNET对整条字段向量赋权，MaskNet可对不同维度分别缩放。连续小权重不会自动跳过后续计算，也不能据此判定字段没有因果作用；门控处理与观测纠偏仍须分别检验。

% Outline: 4.3.1.1.17
\paragraph{Wukong：特征交互的容量扩展}
\label{sec:rs-wukong}

只扩大嵌入表，未必增加组合字段的计算能力。\term{Wukong}以堆叠交互层扩展稠密计算，每层包含因子分解块与线性压缩块\scite{rs-r54}前者从当前向量矩阵$\bm X$计算$\bm X\bm X^\top$等两两交互，将其展开、归一化、经MLP映射并重排为新向量；后者计算$\bm W_{\mathrm L}\bm X$，保留线性组合。两路拼接后与形状匹配的残差相加，再归一化：
\[
 \bm X^{(\ell+1)}=\operatorname{LN}\!\left(
 [\operatorname{FMB}(\bm X^{(\ell)});
  \operatorname{LCB}(\bm X^{(\ell)})]
  +\operatorname{Proj}(\bm X^{(\ell)})\right).
\]
这里方括号沿向量行拼接；形状相同时残差投影可以是恒等。

三字段例子的第一层内积矩阵为
$\left(\begin{smallmatrix}1&0&1\\0&1&1\\1&1&2\end{smallmatrix}\right)$。
因子分解块将这些关系变成新向量，压缩块则保留原字段的线性组合。第二层让两类向量继续相互作用，因此既能沿用低阶信息，也能组合前层已经形成的关系。由于包含MLP、归一化等非线性，交互阶数的直觉不应被误读为整个网络是一条固定阶多项式。

最终响应头用点击等标签训练。论文还利用低秩投影及乘法结合律减少显式内积矩阵的代价；输出向量数、层数、投影秩和MLP宽度共同控制预算。规模收益应限制在论文的数据及资源范围内，不能外推为无限增大必然改善。它扩展的是字段交互容量，字段—历史联合计算和跨阶段共享是另两个问题。

% Outline: 4.3.1.1.18
\paragraph{Hiformer：异质字段的注意力评分}
\label{sec:rs-hiformer}

应用ID、时段和地区虽都表示为向量，却不是语义相同的词元。\term{Hiformer}研究怎样为异质字段安排不同的信息投影\scite{rs-r55}基础异质注意层让字段$a$拥有自己的查询投影，字段$b$拥有自己的键、值投影，并采用字段专属前馈层。于是应用与时段比较时可以突出使用时间，应用与地区比较时可以突出地域可用性。

完整Hiformer进一步使用复合投影：先拼接全部字段，再以一个分块矩阵生成各查询、键和值。例如
\begin{equation}
 [\bm k_1;\ldots;\bm k_F]=
 \widehat{\bm W}_{K}[\bm e_1;\ldots;\bm e_F],
 \label{eq:rs-hiformer}
\end{equation}
其中矩阵第$(a,b)$块允许输出单元$a$的键读取输入字段$b$。它不只是为单个字段各换一个矩阵，而是先生成含跨字段信息的复合单元，再进行注意交互。加入的任务单元从交互层取得表示，由各自MLP头预测点击、安装或观看等目标。

一个标量教学投影可令应用键为$0.8e_{\mathrm{app}}+0.2e_{\mathrm{region}}$，时段键为$e_{\mathrm{hour}}$；相同应用在不同地区便能得到不同键。实际维度下，完整复合矩阵带来约$O(F^2d^2)$的投影成本，可通过低秩分解降低。末层只需任务单元输出时，可以裁剪不再用于预测的查询计算；这不等于前层从未计算过字段。本模型处理字段交互，并没有因此取得逐条长历史的顺序建模能力。

% Outline: 4.3.1.2
\subsubsection{基于历史行为的评分}
\label{sec:rs-history-scoring}

静态字段不能充分说明用户此刻在找什么。第~\ref{sec:rs-user-understanding}~节已经讨论怎样从行为形成表示，评分时还需确定候选何时参与。GRU4Rec、SASRec可先产生用户状态，再与候选表示比较；若迁移到CTR，则需加入候选及情境字段，用曝光点击标签训练响应头。DIEN的兴趣演化已经依赖候选，表示通常随候选变化，不能把它当作一次计算后对所有候选共用的向量。

TransAct将实时序列加工结果交给下游排序，HSTU则可以按其行为转导任务产生候选相关输出。第~\ref{sec:rs-transact}、\ref{sec:rs-hstu}~节的原目标并不自动等于校准CTR。迁移时要交代保留了哪些行为、是否添加响应头、标签窗口如何变化。若一份历史要比较500个候选，独立用户编码与500次候选相关聚合的成本也不能混在一起报告。

% Outline: 4.3.1.2.1
\paragraph{DIN：候选相关兴趣评分}
\label{sec:rs-din-scoring}

同一段登山包、水壶、相机历史，对于背包候选与相机候选应该提供不同证据。DIN让候选查询历史，通过局部激活得到$\bm h(\bm x,i)$，再将其与候选及其他字段拼接，交给CTR预测头\scite{rs-r12}兴趣提取机制见第~\ref{sec:rs-din}~节，这里保留评分接口：
$p_i=\sigma(\operatorname{MLP}[\bm h(\bm x,i);\bm v_i;\bm x_{\mathrm{other}}])$。

假设背包候选使历史中的背包、水壶获得较大权重，相机候选则主要读取相机事件；即使同一预测头处理两者，输入兴趣表示也不同。点击交叉熵会同时更新局部激活、嵌入和预测头。单次请求代价依赖候选数与参与匹配的历史长度，只有不依赖候选的部分能提前复用。

DIN原始应用为广告CTR，推荐复用时需重新定义响应标签。候选相关权重本身不编码历史顺序，也不同于MIND预先形成多个兴趣向量再访问全目录：前者解释给定候选，后者负责取得候选。

\Needspace{85mm}
% Outline: 4.3.1.2.2
\paragraph{SIM：长历史检索与精细评分}
\label{sec:rs-sim-scoring}

若历史跨越多年，每个候选遍历全部行为难以满足预算。SIM让候选先从长历史选出相关子序列，再由精细兴趣模块与其他字段共同产生响应分数\scite{rs-r20}检索和精细匹配的具体机制见第~\ref{sec:rs-sim}~节；本节强调它们必须衔接到同一曝光样本的CTR标签。

例如完整历史含一万条行为，统一最近截断保留100条，可能恰好丢掉旧的摄影兴趣；候选相关检索则可能从较早记录中取回相机、镜头等事件。精细模块只处理选出的$K$条，再接候选响应头。比较时应固定原始历史、候选及标签，同时检查被检索阶段漏掉的证据、最终误差和总时延，不能只计算精细模块的成本。

SIM取回的是历史证据，第~\ref{chap:rs-retrieval}~章取回的是待推荐对象。使用长期历史也不等于优化长期行动价值；后者还要建模本次推荐对后续状态的影响。

% Outline: 4.3.1.3
\subsubsection{基于语义内容的评分}
\label{sec:rs-semantic-scoring}

物品名称、说明和用户表达可以补充ID之外的含义。语义向量可作为字段送入上述评分器，也可让语言模型直接输出评分、接受判断或候选次序。后一种接口要特别检查身份、提示和解析协议：输出了一段合理文字，不等于已经比较了整个目录，更不等于概率经过校准。

% Outline: 4.3.1.3.1
\paragraph{P5：文本化响应与推荐判定}
\label{sec:rs-p5-scoring}

P5将不同推荐任务表达为文本输入和文本目标，用统一的序列到序列模型学习\scite{rs-r21}用于当前候选评分时，输入可以包含任务提示、用户标识、物品标识及可用描述，目标可以是评分数字或肯定、否定判断。其训练目标是目标文本的词元负对数似然，而不是自动对所有输出施加同一个数值回归损失。

例如教学模板为“用户u的已知偏好为……；预测其对物品i的1至5分评分”，监督文本为“4”。服务端应限定合法输出集合，解析为数值后再比较候选；若输出“很好”或超出量表的“8”，应进入明确的无效输出处理，而不能静默当成高分。不同候选必须遵循相同模板和解码协议。

P5还包括序列推荐、直接推荐、解释和评论任务，这些任务具有不同输出语义。生成“是”可以表示一次模型判定，但其词元概率受提示和词表影响，并不自动等于点击概率。用户ID、物品ID及模板切分还可能产生泄漏，须在第~\ref{chap:rs-evaluation}~章的协议下检验。

% Outline: 4.3.1.3.2
\paragraph{TALLRec：语言模型推荐判定}
\label{sec:rs-tallrec}

通用语言模型未必会把历史偏好与指定候选对应起来。\term{TALLRec}将用户历史中的喜欢、不喜欢等证据及目标物品写入推荐指令，用明确的肯定、否定反馈适应模型\scite{rs-r22}参数高效适应只改变训练哪些参数；当前任务仍是“对给定目标作判断”，不是自行访问目录。

教学输入固定一段偏好历史，分别询问通勤背包与摄影镜头。训练为每个目标提供真实任务标签，服务时读取允许的回答，或在约定的肯定、否定词元之间恢复比较分数。若以两个词元概率归一化为
$q=p_{\mathrm{yes}}/(p_{\mathrm{yes}}+p_{\mathrm{no}})$，
应明确这是所选解析协议，还要在目标数据上校准和评价，不能把任意自然语言回答直接换成CTR。

候选标题有时不足以唯一确定物品，内容也可能掩盖用户过去行为的具体对象。应检查身份映射、提示长度、标签数量及领域变化，分别报告低数据量实验的范围。自然语言理由可辅助阅读，但不能证明模型识别了真实偏好原因。

% Outline: 4.3.1.3.3
\paragraph{RankGPT：指令候选重排}
\label{sec:rs-rankgpt}

语言模型也可以直接比较一组文本。\term{RankGPT}把需求和已有候选编号交给模型，要求返回候选排列；长列表则通过重叠滑动窗口逐段重排\scite{rs-s14}原论文主任务是搜索相关性重排，个性化推荐迁移还需提供用户偏好证据，不能把搜索实验当成推荐效果。

假设候选编号为$[1],[2],[3],[4]$，输出$[3]>[1]>[4]>[2]$只表示当前比较的次序。实现必须检查编号是否重复、缺失或越界，规定错误输出的回退，并保留每次窗口输入与恢复后的全局位置。由尾部向前处理的重叠窗口，允许原来靠后的候选逐步向前移动；窗口宽度和重叠步长会改变调用次数及可比较范围。

在同一候选池上改变输入顺序，是检查位置敏感性的一个直接办法。还应把候选覆盖与重排质量分别评价：模型不能把未提供的好物品排到前面。若把重排结果蒸馏给较小评分器，需明确教师次序、训练候选及蒸馏损失；阶段预算在第~\ref{sec:rs-stage-connection}~节讨论。

% Outline: 4.3.2
\subsection{评分模块的组合与统一}
\label{sec:rs-scoring-organization}

字段和历史都能影响分数，但还要决定两者何时相遇。先把历史压成一个向量再做字段交互，计算简单，却可能提前丢失候选需要的细节；逐层交换信息或放入共同骨干，可以保留更多交互机会，也改变复用条件。下面比较完整评分器的组织，不是要求服务依次执行这些模型。

% Outline: 4.3.2.1
\subsubsection{基于分离模块的融合评分}
\label{sec:rs-separate-fusion}

一种直接方案是用DIN得到候选兴趣，拼接用户、候选和情境字段，再交给DCN-V2；也可以让序列编码器先产生用户表示，作为RankMixer的一个输入单元。两者保留历史与字段模块各自的计算结构，但通过共同响应损失联合训练。若历史表示依赖候选，在线仍需按候选重新计算相应部分。

分离模块不必只在末尾拼接。FinalMLP融合两条字段分支，InterFormer让字段和历史模块逐层交换摘要。比较时应明确融合对象、传递信息和发生位置，并固定历史长度、候选数及参数预算；否则更长历史带来的收益容易被归到“融合结构”上。

% Outline: 4.3.2.1.1
\Needspace{90mm}
\paragraph{FinalMLP：互补分支的融合评分}
\label{sec:rs-finalmlp}

两条同构MLP也可以通过不同输入选择形成互补。\term{FinalMLP}用上下文门控分别调整两分支的输入，再以多头双线性聚合两个分支的输出\scite{rs-r53}设字段拼接为$\bm z$，两个门控为$\bm g_1,\bm g_2$，则$\bm h_j=\operatorname{MLP}_j(\bm g_j\odot\bm z)$。融合同时包含两侧线性贡献与双线性作用，单头可以写成
$s=b+\bm w_1^\top\bm h_1+\bm w_2^\top\bm h_2+\bm h_1^\top\bm W\bm h_2$。
多头实现对表示分块，减少完整双线性矩阵的代价。

若两个教学分支输出为$(1,0.5)$和$(0.2,1)$，$\bm W$取单位矩阵，则双线性贡献为0.7；假设偏置与线性项合计为$-1$，最终概率为$\sigma(-0.3)\approx0.426$。点击标签通过同一交叉熵更新门控、两分支及融合层。FiBiNET交互的是字段向量，FinalMLP此处交互的是两个已加工分支的输出，不能只因都用了双线性项便视为同一结构。

% Outline: 4.3.2.1.2
\paragraph{InterFormer：字段与历史的逐层融合评分}
\label{sec:rs-interformer}

只让字段指导历史聚合，没有利用历史反过来修正字段计算。\term{InterFormer}保留字段交互模块与序列模块，让两者在各层通过桥接模块交换摘要\scite{rs-r56}字段摘要进入序列注意和个性化前馈，序列的全局、多方面及近期摘要则进入字段交互。逐条历史和字段表示仍分别保留，不必把两者全部压成同一种单元。

例如当前字段指出候选为通勤背包，第一层字段摘要使序列模块关注通勤事件；序列侧又发现近期多次查看大容量商品，将相应摘要交回字段模块。第二层字段交互便可结合容量、价格与当前兴趣重新计算，而不是一直使用第一层的固定历史摘要。最终字段表示接响应头，以曝光点击交叉熵共同训练两侧模块和桥接。

桥接摘要仍是压缩，粒度和宽度会限制可传递的信息；增加层间交换也会增加计算依赖。它与下文OneTrans的区别在于保留两套模块并交错融合，而不是共用一套逐层骨干。两者都只解决候选评分内部的组织问题。

% Outline: 4.3.2.2
\subsubsection{基于联合骨干的评分}
\label{sec:rs-joint-backbone}

另一种组织方式是让字段和历史持续进入共同的逐层计算。联合骨干并不要求所有参数共享，也不要求任意输入彼此可见。MTGR保留业务交叉并按用户聚合，OneTrans将行为及字段组织为共同序列，MixFormer让字段查询反复交互并读取历史；它们都仍需要外部候选。

\Needspace{110mm}
% Outline: 4.3.2.2.1
\paragraph{MTGR：交叉特征与行为联合评分}
\label{sec:rs-mtgr}

把输入改成序列时，若丢弃已有的用户—候选交叉统计，模型可能同时失去有效证据。\term{MTGR}保留用户字段、历史、即时行为及候选交叉特征，将它们映射到共同维度，以HSTU式编码器产生各候选响应\scite{rs-r40}MTGR的作者释义是Meituan Generative Recommendation，这里的输出是给定候选的分数，并非直接解码整张推荐列表。

模型按用户聚合同一窗口的候选，复用窗口之前的静态历史和用户计算。不同语义组采用相应归一化，动态掩码则控制即时事件及候选之间的可见性。窗口前的静态字段、历史可以被候选读取；窗口内的新行为只对时间更晚的候选可见；候选之间隔离，不能借另一条样本泄漏当前不可用的信息。HSTU基础计算见第~\ref{sec:rs-hstu}~节。

例如上午10点候选$a$、下午2点候选$b,c$被放进同一训练窗口，用户中午12点点击相机。该点击可以供$b,c$评分，不能供$a$评分；$b$之后才得到的响应也不能作为$c$的输入，除非$c$的真实请求时点已经观测到它。各候选仍按自己的曝光标签计算损失。用户计算的复用减少重复工作，但候选表示与输出不会因聚合而完全不花成本；应按整次请求计量。

% Outline: 4.3.2.2.2
\paragraph{OneTrans：特征交互与行为序列联合评分}
\label{sec:rs-onetrans}

历史先汇总成固定向量，会限制后续字段能选择的行为细节。\term{OneTrans}把行为单元放在字段单元之前，由同一Transformer堆栈处理；行为侧共享查询、键、值和前馈参数，不同字段单元则保留各自参数\scite{rs-r41}通用注意计算见\BookChapterRef{chap:transformer}{模型卷的“Transformer”章}，本节关心混合参数和可见范围。

考虑四条行为$b_1,b_2,b_3,b_4$，随后是用户字段$u$和候选字段$i$。因果掩码使$i$读取全部此前行为及$u$，使$u$读取此前行为，但行为不能反向读取后来的候选。于是候选输出已经含有历史和字段的多层交互，可接CTR、CVR等响应头，以各任务对应标签联合训练。多行为输入若按意图而非时间排列，掩码含义也应随其组织方式说明；不能将任意排列都叫时间因果关系。

金字塔计算逐层减少发出查询并保留输出的行为单元。例如某层只保留$b_3,b_4,u,i$的查询，但该层键和值仍来自$b_1$至$i$的完整当前输入；下一层才接收缩短后的输出。这与一开始删除$b_1,b_2$不同：较晚单元已经有机会读取它们。若当前层输入长度为$L$、保留查询数为$L'$，注意乘加约为$O(LL'd)$，前馈只处理保留输出。

历史不能读取候选，使同一请求中历史键值可在多个候选间复用。跨请求复用还依赖追加式历史、相同排序方式及模型版本；事件修改、候选相关子序列或改变位置编码，都可能破坏直接复用条件。新增行为的键值投影可以只做增量部分，但新查询读取旧缓存仍有代价。OneTrans统一评分内部计算，跨阶段联合学习另见第~\ref{sec:rs-onetrans-v2}~节。

% Outline: 4.3.2.2.3
\Needspace{95mm}
\paragraph{MixFormer：特征交互与历史聚合联合评分}
\label{sec:rs-mixformer}

联合评分不一定把所有历史和字段放进同一自注意序列。\term{MixFormer}将非序列字段组成多个查询头，每层依次进行查询混合、历史读取及输出融合\scite{rs-r42}Query Mixer通过头间信息重排及各头前馈交换字段信息，随后这些查询对历史键值作交叉注意，再由各头的前馈计算融合当前字段与历史结果。堆叠后的表示交给多任务响应头。

沿背包候选看两层计算：第一层查询已经包含用户、价格、容量等字段，读取到通勤相关历史；第二层的查询混合继续组合这些结果，新的查询可以强调与容量或近期购买相关的事件。与固定一次历史汇总相比，后层读取条件会变化；与OneTrans相比，历史主要作为被查询的序列，交互位置和参数组织不同。

用户—候选解耦变体进一步限制信息流，使候选侧可以读取用户侧，而用户侧不接收候选信息，从而复用同次请求的用户计算。若允许候选先修改用户查询，就不能再把修改后的结果直接供其他候选使用。响应标签、头数、历史长度和稠密容量需共同记录；联合扩展的收益应连同论文的数据和预算检验，不能仅凭名称推断网络结构或全流程整合范围。

图~\ref{fig:rs-scoring-organization}强调三个不同的融合位置；表~\ref{tab:rs-scoring-reuse}进一步比较可复用部分。共同骨干与共享缓存是相关但不同的选择，后者由信息依赖决定。
\begin{figure*}[htbp]
 \centering
 % 字段与历史的三种融合位置；全版心，约162mm。
\begin{tikzpicture}[
 font=\figurefont\footnotesize,
 box/.style={draw=BookRule,fill=BookPaper,align=center,
   text width=32mm,minimum height=11mm,inner sep=3pt},
 flow/.style={-{Stealth[length=2mm]},draw=BookMuted,line width=.9pt}
]
 \node[font=\figurefont\footnotesize\bfseries,anchor=west] at (-1.8,1.35)
   {末端融合：先独立压缩，再比较候选};
 \node[box,fill=FigureInput!10] (h1) at (0,.55) {原始历史};
 \node[box,fill=FigureInput!10] (f1) at (0,-.85) {用户、候选与情境};
 \node[box,fill=FigureModel!10] (he1) at (4.8,.55) {历史表示};
 \node[box,fill=FigureModel!10] (fi1) at (4.8,-.85) {字段表示};
 \node[box,fill=FigureOutput!10] (o1) at (12,0) {融合及响应头};
 \draw[flow] (h1)--(he1);
 \draw[flow] (f1)--(fi1);
 \draw[flow] (he1.east)--(8,.55)--(o1.west);
 \draw[flow] (fi1.east)--(8,-.85)--(o1.west);
 \node[font=\figurefont\footnotesize\bfseries,anchor=west] at (-1.8,-2)
   {逐层桥接：两套表示通过摘要交换信息};
 \node[box,fill=FigureModel!10] (h2) at (0,-2.85) {历史模块};
 \node[box,fill=FigureModel!10] (f2) at (5.5,-2.85) {字段模块};
 \node[box,fill=FigureOutput!10] (o2) at (12,-2.85) {响应头};
 \draw[flow] (h2.east) to[bend left=20] node[above]{历史摘要} (f2.west);
 \draw[flow] (f2.west) to[bend left=20] node[below]{字段摘要} (h2.east);
 \draw[flow] (f2)--(o2);
 \node[font=\figurefont\footnotesize\bfseries,anchor=west] at (-1.8,-4.3)
   {联合骨干：每层同时安排字段交互与历史读取};
 \node[box,fill=FigureInput!10] (i3) at (0,-5.2) {历史与字段单元};
 \node[box,fill=FigureModel!10] (b3) at (5.5,-5.2) {联合层堆栈\\掩码控制信息可见};
 \node[box,fill=FigureOutput!10] (o3) at (12,-5.2) {候选响应头};
 \draw[flow] (i3)--(b3);
 \draw[flow] (b3)--(o3);
\end{tikzpicture}

 \caption{候选评分中字段与历史相遇的位置。箭头表示信息依赖，图中“联合层”并不表示全部输入双向可见；具体掩码及参数共享仍由各模型规定。}
 \label{fig:rs-scoring-organization}
\end{figure*}

\begin{table*}[htbp]
 \centering\small
 \begin{tabularx}{\linewidth}{@{}lXXXX@{}}
 \toprule
 组织 & 原始输入与交互位置 & 可见范围与参数 & 计算复用 & 输出及主要代价\\
 \midrule
 分离模块融合 & 字段与历史分开加工；末端或逐层摘要融合 & 由模块接口限定；参数分别组织 & 独立用户编码可复用；候选相关聚合须重算 & 响应头；两侧模块及桥接\\
 MTGR & 字段、静态历史、即时行为和候选交叉共同编码 & 动态时间掩码；候选隔离；语义组归一化 & 用户窗口聚合；仍需候选计算 & 多候选响应；窗口大小与掩码\\
 OneTrans & 行为在前、字段在后共同堆叠 & 因果可见；行为共享、字段专属参数 & 请求内历史缓存；跨请求需追加与版本一致 & 响应头；保留查询数与序列长度\\
 MixFormer & 字段查询逐层混合并读取历史 & 查询头参数化；解耦变体限制候选回流 & 解耦后复用请求内用户侧；跨请求另验 & 多任务头；查询数、历史与稠密容量\\
 \bottomrule
 \end{tabularx}
 \caption{评分模块组织与复用条件。所有行均以给定候选为前提；固定信息时点、候选数、标签和预算，才能区分架构变化与更多信息、更多计算带来的收益。}
 \label{tab:rs-scoring-reuse}
\end{table*}

% Outline: 4.3.3
\subsection{观测数据下的评分学习}
\label{sec:rs-observed-learning}

评分器只见到日志中的样本，而日志并非从所有用户—物品对均匀抽取。旧系统决定展示什么，位置影响是否查看，用户点击还可能受流行度信息影响。缺失、误点与恶意反馈又会改变标签含义。处理这些条件时，应先定义事件和可观测信息，再决定重加权、建模或数据治理怎样进入训练。

% Outline: 4.3.3.1
\subsubsection{反馈标签与质量控制}
\label{sec:rs-feedback-quality}

人工相关性判断、用户显式评分、点击、消费和生成样本提供不同监督。人工标注可能只看内容而不知道个体需求；点击要求对象先被展示且被注意到；消费依赖后续追踪；语言模型生成的偏好则是另一个模型的判断。不能把它们统一改成0、1后就认为目标相同。

一份可用样本应能说明请求、候选、展示位置、结果事件及观察截止时点。曝光不等于查看，未点击可能因为没有注意到，点击也可能是误点。若观察窗口为7天，发生1小时的样本仍未成熟；此时“尚未记录购买”与“7天没有购买”不同。延迟回传和广告标签修正见第~\ref{chap:rs-advertising}~章。

质量检查应追踪字段的实际可用时间。最终成交价、事后评分和未来聚合统计不能作为请求时输入；以生成文本补充内容时，也需排除目标标签及未来评论。错误追踪造成的缺失应先修正记录或标记可用范围，不能期待网络自动辨认每种缺失原因。

% Outline: 4.3.3.2
\subsubsection{负例构造与训练采样}
\label{sec:rs-ranking-negatives}

未交互、曝光未点击、人工判为不合适和从目录采样，分别对应不同的负例依据。曝光未点击至少说明进入过展示，但不保证被查看；未交互物品可能从未有机会出现；人工负例的有效范围由标注任务决定。BPR需要偏好比较假设，CTR则需要响应事件及窗口，不能共享一个未经说明的“负标签”。

采样还会改变概率。假设所有正例保留，负例以常数概率$a$独立保留，且除此之外不依赖输入；若真实概率为$p$，采样后概率$q$满足
\[
 q=\frac{p}{p+a(1-p)},\qquad
 p=\frac{aq}{1-q+aq}.
\]
教学例中$a=0.1$、模型在采样数据上预测$q=0.5$，恢复的$p$约为0.0909。输入相关采样时需使用对应的$a(\bm x,i)$；漏掉该修正会把训练分布概率误当成线上概率。

难负例可加强局部竞争，也可能包含真实喜欢但未观察的物品。召回训练面对目录竞争，排序面对上游筛选后的候选；第~\ref{sec:rs-retrieval-training}~节的采样原则可复用，但需要重新说明比较分布。负采样处理的是训练计算与权重，不能替代曝光或位置机制的分析。

% Outline: 4.3.3.3
\subsubsection{曝光与点击选择偏差下的评分学习}
\label{sec:rs-selection-bias}

旧系统只展示过一款背包时，另一款的缺失反馈不能直接当成拒绝。Recommendations as Treatments将反馈是否被观察记为$O$，用观测倾向$o_{u,i}=\mathbb P(O_{u,i}=1\mid\text{已说明的信息})$校正风险\scite{rs-e01}若目标是全体用户—物品对上的平均损失，一种逆倾向估计为
\begin{equation}
 \widehat R_{\mathrm{obs}}
 =\frac1{|U||I|}
   \sum_{(u,i):\,O_{u,i}=1}
    \frac{\ell(y_{u,i},\widehat y_{u,i})}{o_{u,i}}.
 \label{eq:rs-observed-risk}
\end{equation}
观测概率为0.5与0.1的两个样本，其损失权重为2与10；低概率样本代表更多未观察机会，也带来更大方差。

这个估计需要正确的观测机制、足够覆盖及相应可忽略性条件。若某类对象从不被观察，就无法单靠这份日志识别其平均损失；模型外推应与日志支持的范围分开。倾向未知时还需交代估计来源，不能以响应模型自己的点击预测代替观测概率。

广告CVR只在已点击样本中直接取得转化监督，属于进入训练的选择条件，后文用ESMM和ESCM$^2$讨论。观测倾向、位置检查概率、点击概率，以及策略评价中的日志行动概率具有不同事件与分母，必须分别命名。

% Outline: 4.3.3.4
\subsubsection{位置偏差下的评分学习}
\label{sec:rs-position-bias}

即使对象已经展示，低位结果也可能因为较少被查看而缺少点击。令$E$表示查看，$R$表示在当前请求下具有吸引力或相关性，$C$表示点击，$k$为位置。在简化的位置点击模型中，假设点击需要查看及正响应，查看只依赖位置，相关性不再受位置影响，得到
\begin{equation}
 \mathbb P(C=1\mid\bm x,i,k)=e_k r(\bm x,i).
 \label{eq:rs-pbm}
\end{equation}
这个乘积来自生成假设，不是点击的定义。若用户因信任首位而点击原本不喜欢的对象，或者前一结果改变了后续查看行为，简单分解就可能不成立。

图~\ref{fig:rs-position-paths}还表明，旧排序同时影响位置分配，相关性及其代理特征也参与旧排序。因此首位点击率更高，既可能来自更常查看，也可能来自候选本身更相关；直接比较位置均值不能分离二者。随机换位、可比较的历史排序器变更或额外结构假设，才为分解提供证据。
\begin{figure}[htbp]
 \centering
 % 位置、查看与信任的简化关系；正文宽。
\begin{tikzpicture}[
 font=\figurefont\footnotesize,
 box/.style={draw=BookRule,fill=BookPaper,align=center,
   text width=20mm,minimum height=10mm,inner sep=3pt},
 flow/.style={-{Stealth[length=1.8mm]},draw=BookMuted,line width=.9pt}
]
 \node[box,fill=FigureInput!10] (r) at (0,2) {相关性\\及其代理特征};
 \node[box,fill=FigureModel!10] (old) at (3.4,2) {旧排序};
 \node[box,fill=FigureInput!10] (pos) at (6.8,2) {展示位置};
 \node[box,fill=FigureModel!10] (e) at (6.8,0) {查看};
 \node[box,fill=FigureModel!10] (trust) at (3.4,0) {信任};
 \node[box,fill=FigureOutput!10] (c) at (0,0) {点击};
 \draw[flow] (r)--(old);
 \draw[flow] (old)--(pos);
 \draw[flow] (pos)--(e);
 \draw[flow] (r)--(c);
 \draw[flow] (e.south)--(6.8,-1)--(0,-1)--(c.south);
 \draw[flow,dashed] (pos.south west)--(trust.north east);
 \draw[flow,dashed] (trust)--(c);
\end{tikzpicture}

 \caption{位置点击偏差的简化因果关系。实线表示模型中保留的作用，虚线表示可能存在的信任路径；相关性及其可用代理影响旧排序，也影响响应。只观察点击和位置通常不足以唯一分离这些作用。}
 \label{fig:rs-position-paths}
\end{figure}

% Outline: 4.3.3.4.1
\paragraph{Propensity SVM-Rank：基于倾向加权的排序学习}
\label{sec:rs-propensity-svm}

已知检查倾向后，可以让一次低位点击代表更多本来可能被查看的正反馈。\term{Propensity SVM-Rank}对点击对象相对其他候选的排序损失作逆检查倾向加权，再训练带正则的Ranking SVM\scite{rs-s09}若线性分数为$s_i=\bm w^\top\bm\phi(\bm x,i)$，可用成对hinge项$\max(0,1-s_i+s_j)$上界点击对象的名次损失，再乘以$C_i/e_{k_i}$。

假设同一对象在两个位置的检查概率为0.5与0.1，一次点击的权重分别为2与10；若某个比较对的hinge损失为0.4，贡献为0.8与4。较少被查看的对象因此不会仅因低位而在训练中被系统性忽略。这里把点击对象与其他对象比较，是构造相关性加权名次风险，不是断言所有未点击对象都是真实负例。线上只按学到的特征分数排序。

原研究通过少量随机换位标定相对位置倾向；另一项工作从多个历史排序器中寻找同一查询—文档出现在不同位置的记录，提取位置比较证据\scite{rs-s33}后一协议仍依赖位置模型、覆盖和记录可比性，多排序器日志本身不等于随机处理。小倾向会增加方差，截断权重也会改变无偏性，二者都需报告。

% Outline: 4.3.3.4.2
\paragraph{Regression-based EM：基于潜变量估计的倾向学习}
\label{sec:rs-regression-em}

个人搜索中，同一查询—文档可能很少重复，难以逐对统计相关性。\term{Regression-based EM}用特征回归$r(\bm x,i)$共享信息，并把查看与相关性当作点击生成中的潜变量，交替估计其后验和更新参数\scite{rs-s32}

在式~\eqref{eq:rs-pbm}的假设下，点击为1时必有$E=R=1$；点击为0时有多种解释。固定当前估计$e,r$，未点击样本的后验为
\begin{equation}
 \mathbb P(E=1\mid C=0)=\frac{e(1-r)}{1-er},\qquad
 \mathbb P(R=1\mid C=0)=\frac{r(1-e)}{1-er}.
 \label{eq:rs-rem-posterior}
\end{equation}
取教学$e=0.5,r=0.4$，两个后验为0.375与0.25。未点击并没有强迫相关性后验变成0。

E步据当前参数计算这些后验，M步以各位置的查看后验更新检查参数，以相关性后验更新回归。原实现从后验采样相关性标签，并用GBDT学习特征回归；所得位置倾向再供后续加权排序器使用。即使两篇个人文档不重复，它们的特征也可通过回归共享统计依据。

EM提高所设模型的拟合程度，却不能单独证明查看与相关性的分解唯一。模型表达错误、位置覆盖不足及强信任偏差仍会使倾向失准。线上使用最终排序器，不要求实时执行整套EM。

% Outline: 4.3.3.4.3
\paragraph{DLA：基于倾向与相关性互估的排序学习}
\label{sec:rs-dla}

相关性有助于估计查看机会，查看机会又有助于学习相关性。\term{双重学习算法}（dual learning algorithm, DLA）将两者组成联合训练问题\scite{rs-s10}排序模型读取查询—文档特征，位置模型读取位置，分别在列表内形成归一化分数$\widetilde r_i,\widetilde e_i$。以位置1为参考，两类损失可写成
\begin{align}
 \mathcal L_r&=-\sum_{i:C_i=1}
       \frac{\widetilde e_1}{\widetilde e_i}\log\widetilde r_i,\\
 \mathcal L_e&=-\sum_{i:C_i=1}
       \frac{\widetilde r_1}{\widetilde r_i}\log\widetilde e_i.
 \label{eq:rs-dla}
\end{align}
更新一个模型时，另一个模型提供校正权重；训练将两种估计互相联系，而线上只保留排序分数。

假设第三位估计检查分数仅为首位的五分之一，第三位点击在排序损失中就取得5倍权重。另一方面，如果第三位对象本来就被估计为很不相关，位置模型也会据此调整“低点击究竟来自哪里”的判断。列表softmax是训练参数化，不能不加条件地解释成每个对象的独立相关概率。

理论依赖检查、相关性与点击的分解及相应优化条件。两个网络能共同拟合日志，不保证任意数据下的可识别性；若相关性模型缺乏表达能力，或位置还通过信任直接改变点击，互估也可能偏离目标。

% Outline: 4.3.3.4.4
\paragraph{PAL：基于检查与点击分解的响应学习}
\label{sec:rs-pal}

普通CTR模型若把位置作为输入，服务排序时却不知道候选最终位置，只能填默认值。\term{PAL}把位置相关检查概率与查看后点击概率分开：ProbSeen只读取位置，pCTR读取用户、对象及其他特征，以两者乘积拟合日志点击\scite{rs-r92}

设$q_\eta(k)$为检查分支，$r_\theta(\bm x,i)$为响应分支，训练目标为
\[
 \ell_{\mathrm{BCE}}\!\left(C,q_\eta(k)r_\theta(\bm x,i)\right).
\]
教学例中，首位与末位检查概率为0.8、0.2，同一商品查看后点击率为0.25，乘积分别为0.20、0.05。两个分支共同解释位置CTR差异；线上只使用$r_\theta$比较候选，无需提前猜一个位置。

结构成立依赖检查只由位置决定、查看后点击不再随位置变化等假设。若把$q$乘一个常数、$r$除以该常数仍保持合法范围，可能得到相同点击乘积，说明拟合本身不足以固定绝对尺度。原任务是推荐CTR，迁移到搜索时还要重新说明“查看后响应”与相关性标签之间的关系。

% Outline: 4.3.3.4.5
\paragraph{Affine Correction：基于仿射点击模型的排序学习}
\label{sec:rs-affine-correction}

用户可能因信任靠前结果而点击不相关对象，单纯除以检查概率不能去掉这种额外点击。\term{仿射校正}（affine correction）把位置条件下的点击概率写成相关性的缩放加偏移\scite{rs-s34}
\begin{equation}
 \mathbb P(C=1\mid i,k)=\alpha_k r_i+\beta_k,\qquad
 \widehat r_i=\frac{C-\beta_k}{\alpha_k}.
 \label{eq:rs-affine}
\end{equation}
请求条件在式中省略。若查看概率为$e_k$，查看后对相关、无关对象的点击概率分别为$\epsilon_k^+,\epsilon_k^-$，则
$\alpha_k=e_k(\epsilon_k^+-\epsilon_k^-)$，
$\beta_k=e_k\epsilon_k^-$。
需要$\alpha_k>0$及合法概率范围。

假设$\alpha_k=0.5,\beta_k=0.2$，点击时校正贡献为1.6，未点击时为$-0.4$。它们不是概率，但条件期望等于$r_i$，可进入相关性加权的排序目标。纯逆倾向只有缩放，无法消除$\beta_k$这项偏移。

在线使用学到的相关性分数，校正系数只需在训练和评价中按协议使用。系数估计错误、没有跨位置覆盖，或真实信任机制偏离仿射形式，都可能破坏校正。允许负训练贡献也要求使用与目标推导相容的损失，不能把贡献截到$[0,1]$后仍声称保持原无偏性。

% Outline: 4.3.3.4.6
\paragraph{DR-LTR：基于偏好预测与加权残差的排序学习}
\label{sec:rs-dr-ltr}

低位逆倾向权重可能很大，而一般双重稳健估计又要求实际处理可见；点击日志中，用户是否查看通常不可见。Oosterhuis的\term{DR-LTR}用各位置的预期作用代替不可见的实际查看，将偏好预测与点击残差结合\scite{rs-s35}

用式~\eqref{eq:rs-affine}的系数，令日志策略将对象$i$放到位置$k$的边际概率为$\mu(k\mid i)$，定义
$\rho_i=\sum_k\mu(k\mid i)\alpha_k$。
在固定请求的教学记法下，偏好估计具有如下形式：
\begin{equation}
 \widehat r_i^{\mathrm{DR}}
 =\widehat r_i^{\mathrm{reg}}
  +\frac1N\sum_{n=1}^{N}
   \frac{C_{n,i}-\alpha_{k_{n,i}}\widehat r_i^{\mathrm{reg}}
                     -\beta_{k_{n,i}}}
        {\widehat\rho_i}.
 \label{eq:rs-dr-ltr}
\end{equation}
未进入该次展示的对象约定$C=\alpha=\beta=0$。实现可以截断分母，并将对象估计与目标位置权重结合，形成可优化的排序准则。原文还给出与点击模型相容的偏好回归训练。

假设$\widehat r_i^{\mathrm{reg}}=0.4$，当前位置$\alpha=0.5,\beta=0.1$，估计$\widehat\rho_i=0.25$。一次点击的预测点击率为0.3，残差贡献为$(1-0.3)/0.25=2.8$；若只有这条教学记录，基准加残差为3.2。大于1并非计算错误，因为它是校正估计，不是单样本概率。多次记录平均后才用于估计目标。

双重稳健的范围需要完整陈述：位置和信任参数须正确；在此前提下，对每个对象，日志位置分布估计准确且截断不生效，或该对象的偏好回归准确，可以取得无偏估计。回归准确时，残差期望为零；位置分布准确时，残差补回回归误差。它不意味着任意两个网络只需一个正确，更不等同第~\ref{chap:rs-evaluation}~章的一般策略DR。

表~\ref{tab:rs-position-methods}汇总各方法在训练中改变的位置。它们可能在相容假设下组合，年份并不规定必经的替代顺序。
\begin{table*}[htbp]
 \centering\small
 \begin{tabularx}{\linewidth}{@{}lXXXX@{}}
 \toprule
 方法 & 倾向或参数来源 & 生成假设 & 训练校正与线上输出 & 主要失败条件\\
 \midrule
 Propensity SVM-Rank & 换位实验或可比较位置记录 & 查看乘相关性 & 点击名次损失加权；特征分数 & 小倾向、缺覆盖、信任作用\\
 Regression EM & 特征回归与位置参数交替估计 & 查看、相关性潜变量 & 估计倾向供排序训练；最终排序器 & 分解不唯一、回归失配\\
 DLA & 位置与排序模型互估 & 检查与相关性可分解 & 两类加权列表损失；排序分数 & 互估错误、优化与识别条件不足\\
 PAL & 位置分支与响应分支联合拟合 & 位置仅影响查看 & 乘积点击损失；查看后响应 & 绝对尺度、结构失配\\
 仿射校正 & 位置缩放及信任偏移估计 & 仿射点击模型 & 减偏移再缩放；相关性分数 & 系数错误、零缩放或缺覆盖\\
 DR-LTR & 点击模型、日志位置分布及偏好回归 & 正确位置与信任参数 & 基准加残差；训练后排序器 & 超出双重稳健条件、截断偏差\\
 \bottomrule
 \end{tabularx}
 \caption{从带位置点击学习候选评分的方法。相同对象在不同位置的可比记录有助于区分查看和相关性；没有覆盖或额外约束时，结构分支本身不能证明分解正确。}
 \label{tab:rs-position-methods}
\end{table*}

% Outline: 4.3.3.5
\subsubsection{流行度偏差下的评分学习}
\label{sec:rs-popularity-bias}

热门物品可能因为更多曝光而积累点击，也可能因为销量展示影响选择；优质内容及个性化匹配同样可能使其热门。因而“减去热门程度”不能自动恢复偏好。下面两种方法分别改变反事实评分与表示分解，均需保留其结构假设，并与真实干预证据区分。

% Outline: 4.3.3.5.1
\paragraph{MACR：基于反事实评分的流行度纠偏}
\label{sec:rs-macr}

\term{MACR}在用户—物品匹配分支之外，增加用户和物品各自的边际分支，以同一交互标签联合训练主评分和辅助任务\scite{rs-r93}令匹配分数为$m_{u,i}$，两边际分支sigmoid输出为$a_u,b_i$，事实组合为$m_{u,i}a_ub_i$。服务时把匹配分数替换为参考常数$c$，构造反事实组合，再取差：
\begin{equation}
 s_{\mathrm{MACR}}(u,i)=(m_{u,i}-c)a_ub_i.
 \label{eq:rs-macr}
\end{equation}
这里保留了边际分支；差值是候选排序分数，不是校准点击概率。

教学例取$a_u=0.8$，热门候选A的$m=0.6,b=0.9$，候选B的$m=0.8,b=0.4$。事实组合分别为0.432、0.256，A领先；若参考$c=0.5$，校正分数变为0.072、0.096，B领先。参考状态使匹配较弱但边际较强的候选受到更大扣除，同时也说明$c$会影响排序，不能隐藏为无关常数。

论文以因果图解释匹配路径和直接路径，模型可接MF或LightGCN等骨干。但边际分支由交互标签学到什么，仍受曝光、质量和其他未建模因素影响。分支命名及反事实计算并不单独证明真实流行度原因已被识别，效果应按人群、流行度区间及对照条件检查。

% Outline: 4.3.3.5.2
\paragraph{DICE：基于兴趣与从众分离的流行度纠偏}
\label{sec:rs-dice}

\term{DICE}为用户和物品分别学习兴趣、从众两套表示，用流行度关系构造不同辅助监督\scite{rs-r94}最终评分保留两项：
\begin{equation}
 s(u,i)=
 \langle\bm h_u^{\mathrm{int}},\bm v_i^{\mathrm{int}}\rangle+
 \langle\bm h_u^{\mathrm{con}},\bm v_i^{\mathrm{con}}\rangle.
 \label{eq:rs-dice}
\end{equation}
它试图让两个分支解释不同因素，而不是服务时默认删除从众项。

若正例$i$比采样负例$j$更冷门，用户仍选择$i$，在模型假设下，兴趣分支应更偏向$i$，从众分支却应更偏向热门$j$；总点击分数仍应满足$i$领先。若正例更热门，只能为总分和从众提供相应约束，不应强迫兴趣也更偏向正例。两类样本因此产生不同的成对训练信号，再与表示分离损失联合学习。

例如冷门水壶的兴趣、从众得分为1.2、0.1，热门水壶为0.7、0.4，总分分别为1.3、1.1；这个教学分解允许兴趣抵消热门优势。其解释依赖两因素结构、独立性、加性评分及流行度代理。未点击仍可能因为未曝光，数据中的分离效果也不足以证明潜在原因唯一。

% Outline: 4.3.3.6
\subsubsection{隐私约束下的评分学习}
\label{sec:rs-private-scoring}

用户记录不能集中汇集时，评分任务没有消失，但可交换的信息变了。联邦学习规定客户端如何协作更新模型；差分隐私限制输出对单条或单个用户数据的敏感程度。二者职责不同，不能以“只上传梯度”代替隐私分析。这里关注本地图、更新和服务表示分别留在哪里。

% Outline: 4.3.3.6.1
\paragraph{FedPerGNN：分散交互图上的协同评分}
\label{sec:rs-fedpergnn}

\term{FedPerGNN}从用户设备上的交互子图学习用户、物品表示，以已观测评分监督预测；学习服务器聚合受保护更新，可信第三方帮助发现匿名邻居\scite{rs-r122}客户端为伪交互物品加入更新，减弱稀疏更新索引暴露真实交互的风险，再裁剪梯度并添加拉普拉斯噪声。图扩展则让第三方匹配加密物品标识，分发共同交互者的匿名表示，为本地提供高阶关系。

例如设备A和B都评价过电影$i$，原始评分及完整历史留在各自设备。第三方看到可匹配的加密标识，为两侧分发匿名邻居信息；学习服务器收到经保护的模型更新，再聚合为新的共享参数。本地用户表示由本地记录及邻居关系更新。原方案训练后还将用户表示上传到学习服务器以提供个性化服务，因此不能将整条流程表述为“任何用户表示都不离开设备”。

隐私分析需分别覆盖真实交互索引、数值梯度、匿名邻居和最终表示。服务器与第三方的信任分工是条件；多轮更新还会累计隐私消耗。更强噪声、更小裁剪阈值或更多伪物品会影响准确度、通信与计算，不能仅凭原始数据保存在本地就断言没有泄漏。

% Outline: 4.3.3.7
\subsubsection{噪声与恶意反馈下的评分学习}
\label{sec:rs-robust-scoring}

误点、错误标签、虚假交互和恶意客户端更新进入模型的路径不同。参数扰动训练可以增强某种局部稳定性，模型投毒实验则检查攻击者能否改变共享参数。两类工作分别回答防护目标和风险入口，也都不能替代曝光、位置或流行度纠偏。

% Outline: 4.3.3.7.1
\paragraph{APR：参数扰动下的稳健排序训练}
\label{sec:rs-apr}

正负候选分差很小时，小幅参数变化就可能翻转次序。\term{对抗个性化排序}（adversarial personalized ranking, APR）在BPR目标上寻找范数受限的增损参数扰动，再同时降低原始与扰动后的损失\scite{rs-r123}令$\mathcal L_{\mathrm{pair}}$只表示式~\eqref{eq:rs-bpr}的成对数据损失，将参数正则单列，可写成
\begin{align}
 \bm\delta^\star&=\argmax_{\lVert\bm\delta\rVert_2\leq\epsilon}
       \mathcal L_{\mathrm{pair}}(\bm\theta+\bm\delta),\\
 \mathcal L_{\mathrm{APR}}
   &=\mathcal L_{\mathrm{pair}}(\bm\theta)
     +\lambda_{\mathrm{adv}}
       \mathcal L_{\mathrm{pair}}(\bm\theta+\bm\delta^\star)
       +\frac{\lambda}{2}\lVert\bm\theta\rVert_2^2.
 \label{eq:rs-apr}
\end{align}
原文以MF实现对抗矩阵分解，先用BPR初始化，再以梯度方向近似最坏扰动，并交替执行扰动构造和参数更新。

在一维教学例中，用户因子为1，正负物品因子为0.51、0.50，原分差仅0.01。把正物品减0.01、负物品加0.01，总扰动范数约0.0141，就会使分差变成$-0.01$。APR让训练关注这种局部脆弱方向，稳健更新后仍按原评分器比较候选，不在服务时施加恶意扰动。

评价应同时报告普通排序质量、规定扰动下的质量和额外训练成本。参数范数球是明确选择的威胁模型，不等同用户刷量、伪造日志或任意客户端投毒；不能把局部稳健性扩大成对未知攻击的保证。

% Outline: 4.3.3.7.2
\paragraph{FedRecAttack：恶意更新下的评分风险诊断}
\label{sec:rs-fedrecattack}

联邦训练把部分更新交给客户端，也就需要检查上传更新是否可信。\term{FedRecAttack}研究受控客户端怎样利用公开交互近似用户表示，再上传目标导向的恶意梯度，推高指定物品的共享表示和排名\scite{rs-r124}它是模型投毒研究，不是一个服务端防御算法。

设正常客户端的更新改善各自评分误差，受控客户端的更新却集中增加目标商品的得分。服务器聚合后，原本没有对目标商品表现兴趣的用户也可能看到它上升；随后产生的曝光又可能进入新日志。教学分析应沿“公开交互—近似用户表示—恶意更新—共享物品参数—候选次序”检查影响，并明确受控客户端比例、可见信息和更新限制。

只报告目标物品命中率不够，还要看普通用户的排序损失、其他物品受损情况及整体质量，以免把破坏所有排序误当成隐蔽攻击成功。APR扰动的是训练中规定的参数邻域，FedRecAttack改变的是上传更新的完整性；FedPerGNN保护记录或更新隐私，也不自动保证参与者诚实。

% Outline: 4.3.4
\subsection{评分应用与维护}
\label{sec:rs-scoring-maintenance}

模型输出后，还要确定它在什么决策中使用、在哪个阶段执行，以及多久以后仍然有效。概率校准、阶段衔接、时间更新和场景适配可以共同作用于同一评分器。数据撤回则改变允许保留的学习依据，需要清除已有影响，不能只靠加入更新的行为解决。

% Outline: 4.3.4.1
\subsubsection{分数解释、校准与不确定性}
\label{sec:rs-score-calibration}

概率校准要求预测值与同类机会中的响应频率对应。例如所有预测0.1的曝光，长期点击频率应接近0.1。总体平均正确仍可能掩盖页面或人群内的偏差；按分数分箱及关键场景检查，才有助于发现混排和价值计算中的系统误差。单调校准可以保持多数次序，却不能把不同标签定义的分数变成统一效用。

不确定性回答的是“这个预测获得多少证据支持”。两个候选都预测0.1，一个附近有大量成熟样本，另一个来自新内容区域，其可靠程度可能不同。置信区间、模型间分歧、表示密度等各有假设；概率值靠近0或1并不自动意味着模型有充分证据。发生后价值、成本及决策规则还需与概率共同使用。

% Outline: 4.3.4.1.1
\paragraph{PRU：预测相关样本的不确定性估计}
\label{sec:rs-pru}

在点击稀少、正负类别重叠的数据中，靠近任意训练样本未必表示预测可靠。\term{预测相关不确定性}（predictive relevance uncertainty, PRU）先按类别对训练损失拟合高斯混合，选取低损失的预测相关样本，再用这些样本的隐藏表示拟合各类高斯分布\scite{rs-r109}模型训练结合谱归一化与双侧Jacobian正则，鼓励隐藏距离保留输入变化，减少表示塌缩对密度判断的干扰。

设前向计算得到$\bm h$，所选可靠样本形成类别均值$\bm\mu_y$、协方差$\bm\Sigma_y$及混合权重$\omega_y$，则不确定性依据可由
\[
 q(\bm h)=\sum_y\omega_y
   \mathcal N(\bm h;\bm\mu_y,\bm\Sigma_y)
\]
的低密度给出，例如使用负对数密度作为排序量。线上只需一次模型前向及密度计算，无须对每个候选重复采样许多网络。

两个pCTR同为0.1的候选中，A落在预测相关样本的高密度区域，B远离这些区域，可以对B采用更保守回退或在允许范围内安排探索。阈值仍要用验证数据确定。PRU输出不是点击概率，也不是具有覆盖保证的置信区间；高斯假设、类别样本不足和表征失真都会影响它。探索见第~\ref{chap:rs-presentation}~章，集合风险控制见第~\ref{chap:rs-constraints}~章。

% Outline: 4.3.4.2
\subsubsection{召回与排序的阶段衔接}
\label{sec:rs-stage-connection}

召回要从全目录取回对象，粗排要以较小成本筛选较多候选，精排则能给少量候选更多计算。重排还会处理列表关系。假设教学流程把十万候选缩至两千、再缩至一百，前阶段误删的对象通常无法由后阶段恢复；只在已经进入精排的样本上测精排质量，会遗漏这部分损失。

阶段可以共享表示、共同训练或使用教师蒸馏。例如教师对同一候选产生目标分布$q_T$，学生最小化与$q_T$的交叉熵或KL散度，使粗排保留精排认为好的对象。但教师只见过筛选后的候选，蒸馏也会继承这一覆盖限制。语言模型重排放在后阶段时，还需把调用次数、窗口处理和教师输出错误计入总成本。

下面四个方案分别调整用户编码共享、物品复用、跨阶段监督和生成—评分计算。它们仍保留候选与分数接口，不等同直接产生多物品展示列表。

% Outline: 4.3.4.2.1
\paragraph{OneTrans-V2：跨阶段共享与联合学习}
\label{sec:rs-onetrans-v2}

多个阶段分别编码相同历史，会重复消耗计算，也可能学到不一致的用户状态。\term{OneTrans-V2}让召回、粗排和精排读取共享用户骨干，同时保留各自特征单元及输出\scite{rs-r43}召回采用第~\ref{sec:rs-dcgr}~节的DCGR生成物品标识；粗排、精排计算给定候选的响应分数。三个任务共同训练用户表示，精排再向粗排蒸馏，教师logit在蒸馏项中停止梯度。

停止教师梯度只约束蒸馏路径，不表示精排任务不能更新共享骨干。共同任务损失仍由各自样本回传。阶段间和不同曝光的特征单元隔离，避免某次响应或后阶段候选向其他样本泄漏。原OneTrans的字段—历史计算见第~\ref{sec:rs-onetrans}~节。

教学窗口内有请求$t_1,t_2$，中间新增行为$b$。$t_1$的三个阶段只读取前缀$H_{u,t_1}$，$t_2$读取$H_{u,t_1}$加上当时已到达的$b$；不能因两次请求合并训练，就让$t_1$看到$b$。追加式历史和因果掩码允许摊薄重复编码，缓存还要随模型、事件及特征版本维护。共享没有省去阶段执行和候选截断。

评价应同时记录各阶段质量、对精排高价值对象的保留率及总请求成本，区分稀疏专家扩容、参数化与服务优化各自作用。截至2026年10月5日，本工作按2026年9月预印本使用，其实验范围不能被扩大为所有业务的普遍结论。

% Outline: 4.3.4.2.2
\paragraph{IntTower：可复用物品计算的交互粗排}
\label{sec:rs-inttower}

粗排既想保留双塔物品编码可预存的优势，又希望获得比一次点积更细的交互。\term{IntTower}在两塔内使用轻量字段注意，将用户各层表示和物品末层表示投影为多头，再作最大相似度聚合\scite{rs-r60}令用户第$\ell$层第$h$头为$\bm u_{\ell,h}$，物品头为$\bm v_g$，交互分数为
\begin{equation}
 s(u,i)=\sum_{\ell}\sum_{h=1}^{H}
           \max_{1\leq g\leq H}
               \langle\bm u_{\ell,h},\bm v_g\rangle.
 \label{eq:rs-inttower}
\end{equation}
物品各头可以合并存为一个结构化向量，服务时仍需恢复头的分组和最大匹配。

在单层两头教学例中，用户的两个头分别为$(1,0)$与$(0,1)$。候选A的两个头为$(0.8,0.2)$与$(0.1,0.6)$，两次最大值为0.8、0.6，总分1.4。候选B的两个头为$(0.5,0.5)$与$(0.4,0.4)$，总分1.0。计算需要各头间比较，而不是只算拼接向量的一次普通点积；若还有用户层，继续累加对应分数。

将聚合分数经sigmoid变成响应预测后，训练结合点击交叉熵、用户—正物品表示的对比正则及参数正则。物品侧预存、在线用户计算和逐候选交互共同决定时延。模型的原任务是粗排，保留双塔结构不意味着它自动成为普通单向量近邻索引，也不意味着召回、精排已联合训练。

% Outline: 4.3.4.2.3
\paragraph{HCCP：跨阶段一致性与候选覆盖}
\label{sec:rs-hccp}

只模仿精排已曝光样本，粗排就很难学习自己早先滤掉的对象。\term{HCCP}从精排候选、粗排过滤部分、批内物品和目录采样构造不同层次的候选，再组合跨阶段顺序一致性与对比训练\scite{rs-r61}精排序列结合曝光和点击反馈重新组织，全局列表损失对齐较大候选空间，局部列表损失强调曝光空间内的次序；带间隔的对比目标区分不同来源的候选难度。

例如$a,b$进入精排且$a$点击，$c$被粗排截断，$d$来自同批其他用户，$e$来自目录随机采样。$a,b$可以提供曝光响应及局部顺序，$c$有上游筛选信息却没有当前真实点击标签，$d,e$主要提供额外竞争。训练时不能把这五者都解释为同样可靠的已观察正负反馈。

线上仍由粗排模型输出分数并截断。应检查其对精排偏好对象的保留率、被过滤及长尾候选的覆盖，以及后续实际响应。不同负例难度是训练先验，扩大候选空间并不自动证明已去除选择偏差。HCCP改变监督组织，OneTrans-V2改变共享骨干，二者是不同坐标。

% Outline: 4.3.4.2.4
\paragraph{UniR\textsuperscript{2}：生成轨迹与多目标评分共享}
\label{sec:rs-unir2}

生成物品标识与评估给定候选需要的信息并不相同。UniR$^2$把用户条件、标识生成轨迹和候选字段放进一条序列，使用两类查询及不同可见范围\scite{rs-r74}生成查询读取用户前缀和过去代码；评分查询读取画像、完整生成轨迹及候选字段，不再直接读取长历史。候选字段内部可以相互读取，生成路径却不能看到尚未生成的代码。

例如召回依次生成代码$(c_1,c_2,c_3)$，取得候选$i$；评分阶段再补入价格、作者及交叉统计等字段，结合这一生成轨迹形成多目标预测。生成路径使用代码负对数似然，评分路径使用各响应标签的联合损失；目标仍是候选及其响应分数，没有直接确定最终展示列表。

两条路径共享基础注意权重，使用不同前馈网络；评分侧通过LoRA和停止梯度隔离对基础生成计算的更新。可概括为评分投影
$\operatorname{sg}(\bm W\bm x)+\bm B\bm A\bm x$，
其中$\operatorname{sg}$阻断相应基础路径梯度，低秩支路保留评分适应能力。缓存复用的是相同用户条件与生成轨迹，不能据“单次前向”就声称全部参数共同优化或所有信息双向流动。

% Outline: 4.3.4.3
\subsubsection{分布变化下的评分更新}
\label{sec:rs-score-drift}

一次促销会改变价格和候选，一次页面改版会改变查看概率，用户兴趣也可能随时间变化。误差上升后，应该先区分输入更新不及时、标签尚未成熟、条件响应变化及业务目标改变。前两章分别处理需求依据和目录维护，本节关注它们进入评分后是否仍满足训练条件。

如果所有候选概率近似按同一比例偏移，重新校准可能足以修复价值计算；如果用户—物品次序改变，则只调整校准曲线不够。更新可采用滚动窗口、增量学习或重新训练，但需防止短窗口丢掉稳定的低频偏好，也要检查旧阶段输出与新阶段输入是否匹配。

评价不能只看更新后的单日指标。应记录变化发生前后误差、恢复时间、历史能力保留和维护成本，并使用当时可得的信息重放。更快更新还可能更快吸收误点或污染；时间新鲜度本身不是收益保证。

% Outline: 4.3.4.4
\subsubsection{多场景条件下的评分适配}
\label{sec:rs-multiscenario}

同一用户与商品出现在首页、购物车或购后页时，候选上下文和点击基线都可能不同。场景条件描述“在哪里作预测”，响应任务描述“预测什么”；不能把场景数与响应头数混为一谈。相较第~\ref{sec:rs-emcdr}、\ref{sec:rs-conet}~节的跨域迁移，这里主要考虑共享目录下同时存在的页面分布差异。

% Outline: 4.3.4.4.1
\paragraph{STAR：场景共享与专属评分}
\label{sec:rs-star}

完全独立建模浪费共享信息，完全共享又可能掩盖场景差异。\term{STAR}由场景标识选择专属参数，与共享中心参数组合；分区归一化则为不同场景保留相应统计和变换\scite{rs-r57}一层有效权重为$\bm W_{\mathrm c}\odot\bm W_d$，偏置由中心与场景偏置相加，随后作激活；模型还可以通过辅助分支融合场景信息。

用单元教学例看梯度分工：共享权重$w_{\mathrm c}=2$，首页专属权重0.5，购物车专属权重1.5，输入均为0.4且偏置为零，则两个场景logit为0.4、1.2，概率约为0.599、0.769。首页样本更新中心和首页参数，购物车样本更新中心和购物车参数；两个场景通过中心权重交流，专属分支则保留区别。实际网络还会执行场景内归一化及多层变换。

应同时比较整体与场景内误差、校准和资源，防止大场景平均值掩盖小场景退化。该方法需要已知场景和足够样本，不是未知场景或持续漂移的通用解；预测任务之间的共享仍由MMoE等结构另行处理。

% Outline: 4.3.4.4.2
\paragraph{PEPNet：场景与个体先验的条件评分}
\label{sec:rs-pepnet}

场景内部，不同用户、对象和作者也可能需要不同的计算贡献。\term{PEPNet}通过嵌入个性化网络EPNet和参数个性化网络PPNet生成门控，而不是为每个用户独立保存整套网络\scite{rs-r58}门控基本单元采用两层网络，末端以$2\sigma(\cdot)$产生围绕1的缩放。

EPNet将场景先验与停止梯度的通用嵌入拼接，生成逐维门控，再乘到共享嵌入上。PPNet将用户、对象、作者先验与停止梯度的EPNet输出结合，为各任务、各层生成隐藏单元门控。主预测路径仍能更新共享表示，停止梯度阻止的是门控输入这条旁路，并非冻结所有嵌入。

若点击任务某层表示为$(1,0.5)$、门控为$(1.4,0.6)$，缩放后为$(1.4,0.3)$；观看任务门控取$(0.8,1.2)$，相应结果为$(0.8,0.6)$。后续任务头分别拟合自己的响应标签，联合损失更新门控和网络。与MMoE选择专家不同，这里直接调整表示和隐藏单元的贡献。得到多个预测后，最终目标权衡仍由决策规则规定。

% Outline: 4.3.4.5
\subsubsection{数据撤回条件下的评分维护}
\label{sec:rs-unlearning-maintenance}

用户撤回历史后，删除原始日志并不能清除已经进入嵌入、图关系和模型参数的影响。撤回任务要明确删除哪些记录及其派生结果，再取得只依赖允许保留数据的评分器。它与需求变化时吸收新数据不同，也不由差分隐私自动完成。

% Outline: 4.3.4.5.1
\paragraph{RecEraser：分片重训与自适应评分聚合}
\label{sec:rs-receraser}

全量重训是直接参照，但每次撤回都执行可能代价很高。\term{RecEraser}先按用户、物品或交互相似性构造均衡分片，在各片独立训练子模型，再学习注意聚合用户和物品表示\scite{rs-r125}例如用户表示为$\bm h_u=\sum_j a_{u,j}\widetilde{\bm h}_{u,j}$，物品表示也由相应子模型表示聚合；必要投影将子模型空间对齐，再按选定评分函数输出。

若用户撤回一条电影交互，该记录只属于某个固定分片，就从该片删除并重训受影响子模型，同时重训聚合部分。其他独立子模型不必因这一条记录全部重训。分片过多会破坏协同关系，过少则让一次删除涉及较多训练，因此分片数需要按保留质量与删除成本选择。

维护范围不能止于子模型文件。分片若依赖包含该记录的预训练表示，完整流水线与“从未使用该记录重新开始”的参照仍可能不同；邻域、缓存、聚合器及部署后的派生表示也要检查。固定分片框架下的删除性质与完整训练流程的等价性应分别评价，后续第~\ref{chap:rs-evaluation}~章进一步讨论用户删除及协同影响。

% Outline: 4.3.5
\subsection{评分效果评价}
\label{sec:rs-scoring-evaluation}

响应概率、数值价值和排序次序各需要相应检查。点击概率可用交叉熵、Brier分数及分组校准，数值价值可用MAE、RMSE及关键区间误差；AUC衡量正负区分，NDCG等截断指标检查头部次序。AUC不检验概率尺度，低平均回归误差也不保证前几名正确。增量排序则需处理与对照证据，不能直接以普通点击NDCG代替。

这些指标都受标签窗口、采样和观测机制影响。比较评分模型时，固定请求时信息、候选和预算；比较阶段方案时，记录候选截断、教师覆盖及全流程成本。要区分新增历史、更多字段、更多计算和更及时更新各自的贡献，消融比较才有可解释的对象。

即使以上检查通过，评分器仍只是展示决策的一部分。最终列表的重复、位置、候选之间的竞争与多方约束，可能使单对象预测改善没有转化为实际收益。第~\ref{chap:rs-evaluation}~章建立完整离线、策略和实验协议，下一章则研究怎样把这些分数组织成用户实际看到的展示。

% Outline: 4.4
\section{本章小结}
\label{sec:rs-ranking-summary}

给定候选后的比较必须从目标量开始。响应预估描述在明确条件下发生什么，价值预估还要结合结果的价值；偏好排序学习相对次序，增量排序则比较推荐与不推荐造成的结果差异。相同评分架构可以服务这些任务，但标签、损失和识别证据不能随名称一起互换。

字段交互、候选相关历史与语义内容提供评分依据，分离模块或联合骨干决定这些依据何时相遇。评分内部联合计算和跨阶段共享具有不同范围，复用又取决于时间与信息可见性。可靠学习还需要理解观测选择、位置及流行度作用，明确隐私保护与恶意更新的职责，并在数据撤回时处理已经学入的影响。

概率校准、证据可靠程度、排序质量和最终收益需要分别检验。即使每个候选都预测准确，也不能据此确定它们应该怎样并列、排位或反复出现。下一章从单对象分数转向展示列表与交互，将“谁更合适”进一步落实为“用户这一次会看到什么”。
